% =====================================================================
%  U16：PDF p.347--p.374（页脚印刷 p.337--p.364）
%  §10.4 The tamely ramified case（10.4.1--10.4.6）
%  本片段从节中部开始，不写 \chapter；节编号由 main.tex 中的第 10 章接续。
%  注意：PDF p.347 上半页仍属 §10.3（猜想 10.3.12--10.3.13 与 (10.3-29)），
%  不属于本单元（见 reports/U16.md 的衔接说明）。
%  另：PDF p.374 末句在 PDF p.375 续完，由 U17 单元接译。
% =====================================================================

\section{温和分歧情形}

\begin{keypoints}
    \item \textbf{核心思想：} 把“含 Iwahori 不动向量的表示”范畴化：对应的 Langlands 参数是温和分歧的 ${}^L G$-局部系（正则奇点 + 幂幺单值性），而对应的范畴是固定中心特征 $\chi$ 的 $(\hat{\fg}_{\kappa_c},I)$ Harish--Chandra 模范畴。
    \item \textbf{数学工具：} Kazhdan--Lusztig 对仿射 Hecke 代数 $H(G(F),I)$ 的 $K$-理论实现；Springer 纤维与 $K$-理论；非约化（“加厚”）Springer 分解与 DG 纤维；Miura 变换与 Wakimoto 模。
    \item \textbf{关键结论：} 主猜想 (10.4-13) 断言 $D^b(\hat{\fg}_{\kappa_c}\text{-}\mathrm{mod}_{\mathrm{nilp}})^{I^0}\simeq D^b(\mathrm{QCoh}(\widetilde{\mathrm{MOp}}{}^0_{{}^L G}))$，它是温和分歧局部 Langlands 对应的范畴类比。
    \item \textbf{记号提醒：} $I^0=[I,I]$；$W_{\overline{\nabla}}$ 是 Wakimoto 模；$\mu$ 是 Miura 变换；$\kappa_c$ 为临界水平；$L$-包由三元组 $(\gamma,u,\rho)$ 参数化。
\end{keypoints}

% [Tutor Note: p.337 节引言。本节把上一节的无分歧情形推广到 Iwahori 不动向量：无分歧对应 $K_0=G[[t]]$ 的不动向量，温和分歧对应 $I$ 的不动向量。几何一侧的关键词是“正则奇点 + 幂幺单值性”，它们对应 opers 的留数（同 $q=1$ 时的 $\gamma u\gamma^{-1}=qu$）。阅读时请把 10.4.1（经典分类）与 10.4.2--10.4.3（范畴猜想）分开看。]
上一节我们考察了局部非阿基米德域 $F$ 上约化群 $G(F)$ 的不可约无分歧表示的范畴类比物。回忆这些表示是含极大紧子群 $K_0\subset G(F)$ 的非零不动向量的表示。相应的 Langlands 参数是从 Weil--Deligne 群 $W'_F$ 到 ${}^L G$ 的无分歧容许同态，即经过商
\[
    W'_F\longrightarrow W_F\longrightarrow\hat\Z
\]
分解、且在 ${}^L G$ 中的像为半单的同态。这样的同态由 ${}^L G$ 中的半单共轭类参数化。

我们已经看到，$G(F)$ 的无分歧表示的范畴类比物是范畴 $\hat{\fg}_{\kappa_c}\text{-}\mathrm{mod}_\chi$（带圈群 $G((t))$ 的作用），其中 $\chi$ 是 $D^\times$ 上的 ${}^L G$-oper，其底层的 ${}^L G$-局部系平凡。这些范畴可以称为无分歧的，意为它们含非零的 $G[[t]]$-等变对象。相应的 Langlands 参数是 $D^\times$ 上的平凡 ${}^L G$-局部系 $\sigma_0$，它应当看作无分歧同态 $W'_F\to{}^L G$ 的类比物。然而局部系 $\sigma_0$ 可由许多不同的 opers 实现，这给我们的图像带来额外的复杂性：最终我们需要证明上述类型的 opers 所对应的范畴 $\hat{\fg}_{\kappa_c}\text{-}\mathrm{mod}_\chi$ 彼此等价。特别地，猜想 10.3.10 描述了当 $\chi\in\Op_{{}^L G}(D^\times)$ 时我们期待发生什么。\transnote{这里把“无分歧”的极大紧子群 $K_0=G[[t]]$ 换成 Iwahori 子群 $I$：$I$ 是 Borel 子群 $B$ 在求值映射 $G[[t]]\to G$ 下的原像，故 $I$-不动向量只“看见”惯性群的温和分歧部分；Langlands 参数也随之从“经 $W_F\to\hat\Z$ 分解的无分歧同态”放宽为“经 $\Z\ltimes\C$ 分解的温和分歧同态”。这是 §10.4 全部构造的出发点，也说明了为什么几何侧参数仍是 $D^\times$ 上的 ${}^L G$-局部系。[来源：本书 p.337–338；[Kazhdan and Lusztig 1987]（本书 p.338 引用）。]}

下一步自然是考虑 $G(F)$ 中含 Iwahori 子群 $I\subset G[[t]]$ 的不动向量的表示的范畴类比物。这里 $I$ 是固定 Borel 子群 $B\subset G$ 在求值同态 $G[[t]]\to G$ 下的原像。我们以回顾 D. Kazhdan 与 G. Lusztig [1987] 以及 V. Ginzburg [1997] 关于这些表示的分类来开始本节。然后按 [Frenkel and Gaitsgory 2006c; Frenkel and Gaitsgory 2006a; Frenkel and Gaitsgory 2005; Frenkel and Gaitsgory 2006b] 讨论这些表示的范畴类比物，以及经典图像与几何图像之间的精细互动。本节最后给出 $\fg=\fsl_2$ 情形的一些显式计算，用来说明一般理论。

\subsection{温和分歧表示}

\begin{keypoints}
    \item \textbf{核心思想：} 含 $I$-不动向量的不可约表示对应温和分歧同态 $W'_F\to{}^L G$；当 $G\neq GL_n$ 时出现两个新现象：$L$-包与“并非 $L$-包成员都含 $I$-不动向量”。
    \item \textbf{数学工具：} 三元组 $(\gamma,u,\rho)$ 参数化，$\gamma u\gamma^{-1}=qu$；Springer 纤维 $\operatorname{Sp}_u$ 及其 $\gamma$-不动点的同调；仿射 Hecke 代数 $H(G(F),I)$ 的 $K$-理论实现。
    \item \textbf{关键结论：} 温和局部 Langlands 对应 (10.4-6) 给出三元组与含 $I$-不动向量的不可约表示之间的双射（Kazhdan--Lusztig）。
    \item \textbf{记号提醒：} $W'_F=W_F\ltimes\C$；$C(\gamma,u)$ 是 $\gamma$ 与 $u$ 在 ${}^L G$ 中共同中心化子的分量群；$H_\bullet(\operatorname{Sp}_u^\gamma)$ 是 $\gamma$-不动子簇的同调。
\end{keypoints}

% [Tutor Note: p.338。本节第一个要点：$I$-不动向量对应“温和分歧”参数，因为 Iwahori 子群 $I$ 是 $K_0=G[[t]]$ 中 Borel 的原像（它含第一同余子群 $K_1$ 作为正规子群）；$I$-不动向量只“看见”惯性群的温和部分。第二个要点：$GL_n$ 情形（Bernstein--Zelevinsky）没有 $L$-包问题，因为中心化子连通；对其他约化群，$L$-包与分量群 $C(\gamma,u)$ 必须一并考虑。]
对应于 $G(F)$ 的含 $I$-不动向量的不可约表示的 Langlands 参数，是温和分歧同态 $W'_F\to{}^L G$。回忆 1.1.3 节：$W'_F=W_F\ltimes\C$。称同态 $W'_F\to{}^L G$ 为温和分歧的，若它经过商
\[
    W'_F\longrightarrow \Z\ltimes\C
\]
分解。根据关系式 (1.1-1)，群 $\Z\ltimes\C$ 由两个元素 $F=1\in\Z$（Frobenius）与 $M=1\in\C$（单值性）生成，它们满足关系
\begin{equation}
    FMF^{-1}=qM.
    \tag{10.4-1}
\end{equation}
在容许的温和分歧同态下，生成元 $F$ 映为半单元素 $\gamma\in{}^L G$，生成元 $M$ 映为幂幺元素 $N\in{}^L G$。由公式 (10.4-1)，它们须满足关系
\begin{equation}
    \gamma N\gamma^{-1}=N^q.
    \tag{10.4-2}
\end{equation}
也可以写 $N=\exp(u)$，其中 $u$ 是 ${}^L\fg$ 的幂零元。这时上述关系变成
\[
    \gamma u\gamma^{-1}=qu.
\]
于是，在以下两类对象的等价类集合之间有双射：\transnote{(10.4-1)–(10.4-3) 是 Deligne–Langlands 参数的标准形状：$F$ 的像半单、$M$ 的像幂幺，$q$ 是 $F$ 的剩余域大小；“容许”还要求这两部分在 ${}^L G$ 中的像满足 (10.4-1) 的换位关系。令 $M=\exp u$ 把幂幺元换成幂零元，则条件线性化为 $\gamma u\gamma^{-1}=qu$——正是后文把“留数幂零”与 Springer 纤维对接的形式。注意这里 $u$ 是 ${}^L\fg$ 的幂零元，而不是群元素。[来源：本书 p.338–339；[Lusztig 1995]（本书 p.340 引用）。]}
\begin{equation}
    \mbox{\fbox{\parbox[c]{4.6cm}{\centering 温和分歧容许同态\\ $W'_F\to{}^L G$}}}
    \quad\Longleftrightarrow\quad
    \mbox{\fbox{\parbox[c]{5.2cm}{\centering 对 $(\gamma,u)$：$\gamma\in{}^L G$ 半单，\\ $u\in{}^L\fg$ 幂零，$\gamma u\gamma^{-1}=qu$}}}
    \tag{10.4-3}
\end{equation}
两种情形下的等价关系都是 ${}^L G$ 中元素的共轭作用。

% [Tutor Note: p.339。(10.4-3) 的关键是“容许”（admissible）一词：它要求 $F$ 的像半单、$M$ 的像幂幺，并且 Jordan 分解与 $W'_F$ 的结构相容。若只要求 $M$ 幂幺而不要求 $\gamma u\gamma^{-1}=qu$，就会漏掉与 Frobenius 本征值的匹配；这个方程是 (10.4-1) 在 $\exp$ 下的形式，注意 $q$ 是剩余域大小而不是任意参数。]
现在，对每个这样的 Langlands 参数，我们希望配上 $G(F)$ 的一个含非零 $I$-不动向量的不可约光滑表示。结果表明，若 $G=GL_n$，[Bernstein and Zelevinsky 1977] 证明了以下两类对象的等价类集合之间确实有双射：
\begin{equation}
    \mbox{\fbox{\parbox[c]{4.6cm}{\centering 温和分歧容许同态\\ $W'_F\to GL_n$}}}
    \quad\Longleftrightarrow\quad
    \mbox{\fbox{\parbox[c]{5.0cm}{\centering $GL_n(F)$ 的不可约表示\\ $(R,\pi)$，$R^{\pi(I)}\neq0$}}}
    \tag{10.4-4}
\end{equation}
然而对其他约化群，这样的双射不再成立：出现了两个新现象，下面就来讨论。

第一个新现象是 $L$-包的出现。人们不再期待能把单个不可约光滑表示配给一个特定的容许同态 $W'_F\to{}^L G$。取而代之的是配上一有限族这样的表示（更精确地说，一有限族表示的等价类），称为一个 $L$-包。为区分给定 $L$-包中的表示，需要引入一个附加参数。下面很快会看到在手头的情形中如何做到。然而——这就是上面提到的第二个微妙之处——并非与给定温和分歧同态 $W'_F\to{}^L G$ 相联系的 $L$-包中的所有不可约表示都含非零 $I$-不动向量。幸运的是，用来区分 $L$-包中表示的那个附加参数有某个性质，能告诉我们相应表示是否含 $I$-不动向量。

在温和分歧同态 $W'_F\to{}^L G$ 的情形，这个附加参数是有限群 $C(\gamma,u)$ 的不可约表示 $\rho$；$C(\gamma,u)$ 是 $\gamma$ 与 $u$ 在 ${}^L G$ 中共同中心化子的分量群，并且要求 ${}^L G$ 的中心在其上平凡作用（见 [Lusztig 1995]）。在 $G=GL_n$ 的情形，这些中心化子总是连通的，所以这个参数从不出现。但对其他约化群 $G$，这个分量群常常非平凡。最简单的例子是 ${}^L G=G_2$ 而 $u$ 是李代数 ${}^L\fg$ 的次主幂零元。\footnote{“次主”（subprincipal）意为该元素的伴随轨道在幂零锥中的余维数为 $2$。} 此时对某些满足 $\gamma u\gamma^{-1}=qu$ 的 $\gamma$，分量群 $C(\gamma,u)$ 是对称群 $S_3$，它有三个不可约表示（在同构意义下）。它们各自对应与由 $(\gamma,u)$ 定义的温和分歧同态 $W'_F\to{}^L G$ 相联系的 $L$-包中的一个成员。于是这个 $L$-包由 $G(F)$ 的三个（等价类意义下的）不可约光滑表示组成。然而并非它们都含非零 $I$-不动向量。

% [Tutor Note: p.340 易错处。$L$-包的规模由 $C(\gamma,u)$ 的不可约表示个数决定（此处 $|S_3|=6$ 给出 3 个不可约表示，故 $L$-包有 3 个成员）。但只有那些在 $H_\bullet(\operatorname{Sp}_u^\gamma)$ 中出现的 $\rho$ 才对应含 $I$-不动向量的表示——这是 Lusztig 的“Springer 对应”判据，不是自动成立的。用 $G_2$ 的次主幂零元做例子正因为它给出 $S_3$ 与符号表示的缺失。]
对应于含 $I$-不动向量的 $G(F)$ 表示的 $C(\gamma,u)$ 的表示 $\rho$，由以下性质刻画。考虑 Springer 纤维 $\operatorname{Sp}_u$。回忆（见公式 (9.3-10)）
\begin{equation}
    \operatorname{Sp}_u=\{\mathfrak{b}'\in{}^L G/{}^L B\mid u\in\mathfrak{b}'\}.
    \tag{10.4-5}
\end{equation}
群 $C(\gamma,u)$ 作用在 $\operatorname{Sp}_u$ 的 $\gamma$-不动点子簇 $\operatorname{Sp}_u^\gamma$ 的同调上。$C(\gamma,u)$ 的表示 $\rho$ 对应含非零 $I$-不动向量的 $G(F)$ 表示，当且仅当 $\rho$ 出现在 $\operatorname{Sp}_u^\gamma$ 的同调 $H_\bullet(\operatorname{Sp}_u^\gamma)$ 中。

在 $G_2$ 的情形，次主元素 $u$ 的 Springer 纤维 $\operatorname{Sp}_u$ 是四条射影直线之并，它们按 $D_4$ 的 Dynkin 图的方式彼此连接。对某些 $\gamma$，集合 $\operatorname{Sp}_u^\gamma$ 是一条射影直线（对应 $D_4$ 的 Dynkin 图中心顶点）与三个点（各自位于其余三条射影直线上）之并。相应的群 $C(\gamma,u)=S_3$ 在 $\operatorname{Sp}_u^\gamma$ 上平凡地作用于那条射影直线，并置换那三个点。因此 $S_3$ 的平凡表示与二维表示出现在 $H_\bullet(\operatorname{Sp}_u^\gamma)$ 中，而符号表示不出现。与前两个表示对应的 $G(F)$ 的不可约表示含非零 $I$-不动向量，而对应 $S_3$ 的符号表示的那个不含。\transnote{这两段解释 $L$-包为何“不整体出现”：$L$-包的大小由 $C(\gamma,u)$ 的不可约表示数决定（$S_3$ 有三个不可约表示，故 $L$-包有三个成员），而“是否含 $I$-不动向量”由 Lusztig 的 Springer 判据另行筛选——$\rho$ 必须出现在 $\gamma$-不动 Springer 纤维的同调 $H_\bullet(\operatorname{Sp}_u^\gamma)$ 中。$G_2$ 的次主幂零元给出最干净的例子：$\operatorname{Sp}_u$ 是四条 $\PP^1$ 按 $D_4$ 图连接，$\operatorname{Sp}_u^\gamma$ 为一条 $\PP^1$ 加三个点，故 $S_3$ 的平凡表示与二维表示出现、符号表示缺失。这个判据正是后文把 $\Rep C(\gamma,u)$ 的因子从模范畴一侧翻译到几何一侧的原型。[来源：本书 p.340–341；[Lusztig 1995]、[Chriss and Ginzburg 1997]（本书 p.340–341 引用）。]}

含 $I$-不动向量的 $G(F)$ 表示的局部 Langlands 对应的最终形式如下（这里和 [Kazhdan and Lusztig 1987; Chriss and Ginzburg 1997] 中一样，假定群 $G$ 分裂且中心连通）：
\begin{equation}
    \mbox{\fbox{\parbox[c]{5.6cm}{\centering 三元组 $(\gamma,u,\rho)$，$\gamma u\gamma^{-1}=qu$，\\ $\rho\in\Rep C(\gamma,u)$ 出现在 $H_\bullet(\operatorname{Sp}_u^\gamma,\C)$ 中}}}
    \Longleftrightarrow
    \mbox{\fbox{\parbox[c]{4.8cm}{\centering $G(F)$ 的不可约表示\\ $(R,\pi)$，$R^{\pi(I)}\neq0$}}}
    \tag{10.4-6}
\end{equation}
同样，这应理解为所列对象的两类等价类集合之间的双射。这个双射属于 [Kazhdan and Lusztig 1987]（另见 [Chriss and Ginzburg 1997]）。它由 Deligne 与 Langlands 猜想，随后 Lusztig 作了修改（即加入 $\rho$）。

% [Tutor Note: p.340 背景。tame local Langlands correspondence（Deligne--Langlands 猜想）的证明路径是：把 $G(F)$ 的表示换成仿射 Hecke 代数 $H(G(F),I)$ 的模（A. Borel 定理），再用 $H(G(F),I)$ 的 $K$-理论实现把模编码为 Springer 纤维的等变 $K$-理论。这个“绕道代数”的思路正是下一小节范畴化的模板。]
如何建立这个双射？想法是把 (10.4-6) 右端的 $G(F)$ 的不可约表示换成相应 Hecke 代数 $H(G(F),I)$ 的不可约模。回忆 10.2.1 节：这是 $G(F)$ 上紧支集 $I$ 双不变函数关于卷积构成的代数。它自然作用在任何 $G(F)$ 光滑表示的 $I$-不变向量空间上（见公式 (10.2-2)）。于是我们得到从 $G(F)$ 的光滑表示范畴到 $H(G(F),I)$-模范畴的函子。

根据 A. Borel 的一个定理 [1976]，它诱导了 $G(F)$ 的含非零 $I$-不变向量的不可约光滑表示的等价类集合与不可约 $H(G(F),I)$-模的等价类集合之间的双射。\transnote{这条路线是本节范畴化的模板：先用 Borel 定理把 $G(F)$ 的含 $I$-不动向量的不可约表示换成 $H(G(F),I)$ 的不可约模；再用 Kazhdan–Lusztig 同构 (10.4-8) 把 $H(G(F),I)$ 实现为 Steinberg 簇的等变 $K$-理论；最后在中心特征标 $(\gamma,q)$ 处特殊化，并用 $\gamma$-不动点同调分解出 $C(\gamma,u)$ 的各分量。后文 (10.4-13) 的导出范畴等价正是把这条链上的每个箭头范畴化。[来源：本书 p.341–342；A. Borel, Admissible representations of a semi-simple group over a local field with vectors fixed under an Iwahori subgroup, Invent. Math. 35 (1976)（本书 p.341 引用）。]}

代数 $H(G(F),I)$ 称为仿射 Hecke 代数，它有生成元与关系的标准描述。但为我们的目的，需要 [Kazhdan and Lusztig 1987; Chriss and Ginzburg 1997] 给出的另一种描述，它把该代数等同于 Steinberg 簇的等变 $K$-理论：
\[
    \mathrm{St}=\widetilde{\mathcal{N}}\times_{\mathcal{N}}\widetilde{\mathcal{N}},
\]
其中 $\mathcal{N}\subset{}^L\fg$ 是幂零锥，$\widetilde{\mathcal{N}}$ 是 Springer 分解
\[
    \widetilde{\mathcal{N}}=\{x\in\mathcal{N},\ \mathfrak{b}'\in{}^L G/{}^L B\mid x\in\mathfrak{b}'\}.
\]
于是 $\mathrm{St}$ 的点是由 ${}^L\fg$ 的一个幂零元素与含该元素的两个 Borel 子代数组成的三元组。群 ${}^L G\times\C^\times$ 自然作用在 $\mathrm{St}$ 上：${}^L G$ 共轭三元组的成员，$\C^\times$ 通过乘法作用在幂零元素上，
\begin{equation}
    a\cdot(x,\mathfrak{b}',\mathfrak{b}'')=(a^{-1}x,\mathfrak{b}',\mathfrak{b}'').
    \tag{10.4-7}
\end{equation}
根据 [Kazhdan and Lusztig 1987; Chriss and Ginzburg 1997] 的一个定理，有同构
\begin{equation}
    H(G(F),I)\simeq K^{{}^L G\times\C^\times}(\mathrm{St}).
    \tag{10.4-8}
\end{equation}
右端是 $\mathrm{St}$ 的 ${}^L G\times\C^\times$-等变 $K$-理论。它关于自然的卷积运算构成代数（细节见 [Chriss and Ginzburg 1997]）。它还是其中心上的自由模，中心同构于
\[
    K^{{}^L G\times\C^\times}(\mathrm{pt})=\Rep{}^L G\otimes_\C\C[\mathbf{q},\mathbf{q}^{-1}].
\]
在同构 (10.4-8) 下，元素 $\mathbf{q}$ 映为仿射 Hecke 代数 $H(G(F),I)$ 的标准参数 $q$（这里把 $H(G(F),I)$ 看作 $\C[\mathbf{q},\mathbf{q}^{-1}]$-模）。

现在，代数 $K^{{}^L G\times\C^\times}(\mathrm{St})$、从而代数 $H(G(F),I)$ 有一族自然的模，它们恰由上面那样的对 $(\gamma,u)$ 的共轭类参数化。在这些模上 $H(G(F),I)$ 通过中心特征标作用，该特征标对应于 $\Spec(\Rep{}^L G\otimes_\C\C[\mathbf{q},\mathbf{q}^{-1}])$ 中的一点，即一对 $(\gamma,q)$，其中 $\gamma$ 是 ${}^L G$ 中的半单共轭类，$q\in\C^\times$。在我们的情形 $q$ 是 $F$ 的剩余域的元素个数（从而是素数的幂），但下文将允许 $q$ 取更大的范围：所有非零复数，除去单位根。\transnote{允许 $q$ 取非单位根复数，是为了让中心 $\Rep{}^L G\otimes\C[\mathbf q,\mathbf q^{-1}]$ 的局部化不出现额外退化（$q^n=1$ 时会有额外关系）；几何情形相当于把这个“代数参数”形式地取为 $1$，此时 Hecke 代数的 $q$ 消失，只剩 D-模的 Hecke 作用。这一“先把 $q$ 一般化、再取极限”的操作，是无分歧/温和分歧情形通向整体几何 Langlands 的关键技术步骤。[来源：本书 p.342；[Kazhdan and Lusztig 1987]（本书 p.341–342 引用）。]}

% [Tutor Note: p.341--342 易错处。$K^{{}^L G\times\C^\times}(\mathrm{St})$ 的中心是 $\Rep{}^L G\otimes\C[\mathbf{q},\mathbf{q}^{-1}]$；半单参数 $\gamma$ 与 $\mathbf{q}$ 一起给出中心特征标。当 $q$ 不是单位根时，$\C[\mathbf{q},\mathbf{q}^{-1}]$ 的局部化才有意义（避免 $q^n=1$ 时的退化）。“把 $q$ 一般化”是几何化的关键步骤：经典情形 $q$ 被剩余域固定，几何情形允许 $q$ 变动，从而可以取 $q=1$ 作为极限。]
考虑 $H(G(F),I)$ 模去由 $(\gamma,u)$ 定义的中心特征标所得的商。这恰是代数 $K^{{}^L G\times\C^\times}(\mathrm{St})$ 在 $(\gamma,q)$ 处特殊化，记作 $K^{{}^L G\times\C^\times}(\mathrm{St})_{(\gamma,q)}$。

现在对幂零元 $u\in\mathcal{N}$ 考虑 Springer 纤维 $\operatorname{Sp}_u$。条件 $\gamma u\gamma^{-1}=qu$ 意味着 $u$、从而 $\operatorname{Sp}_u$，在 $(\gamma,q)\in{}^L G\times\C^\times$ 的作用下稳定（见公式 (10.4-7)）。设 $A$ 是 ${}^L G\times\C^\times$ 中含 $(\gamma,q)$ 的最小子代数群。代数 $K^{{}^L G\times\C^\times}(\mathrm{St})_{(\gamma,q)}$ 自然作用在 $(\gamma,q)$ 处特殊化的等变 $K$-理论 $K^A(\operatorname{Sp}_u)$ 上：
\[
    K^A(\operatorname{Sp}_u)_{(\gamma,q)}=K^A(\operatorname{Sp}_u)\otimes_{\Rep A}\C_{(\gamma,q)}.
\]
已知 $K^A(\operatorname{Sp}_u)_{(\gamma,q)}$ 同构于 $\operatorname{Sp}_u$ 的 $\gamma$-不动子集 $\operatorname{Sp}_u^\gamma$ 的同调 $H_\bullet(\operatorname{Sp}_u^\gamma)$（见 [Kazhdan and Lusztig 1987; Chriss and Ginzburg 1997]）。于是得到：$K^A(\operatorname{Sp}_u)_{(\gamma,q)}$ 是 $H(G(F),I)$ 上的模。

不幸的是，这些 $H(G(F),I)$-模一般不可约，要描述 $H(G(F),I)$ 的不可约模还需更多工作。对 $G=GL_n$ 可以证明这些模每个都有唯一的不可约商，这样就恢复了双射 (10.4-4)。但对一般的群 $G$，有限群 $C(\gamma,u)$ 开始起作用。即：群 $C(\gamma,u)$ 作用在 $K^A(\operatorname{Sp}_u)_{(\gamma,q)}$ 上，且这个作用与 $K^{{}^L G\times\C^\times}(\mathrm{St})_{(\gamma,q)}$ 的作用交换。因此有 $K^A(\operatorname{Sp}_u)_{(\gamma,q)}$ 作为 $C(\gamma,u)\times H(G(F),I)$ 的表示的分解
\[
    K^A(\operatorname{Sp}_u)_{(\gamma,q)}
    =\bigoplus_{\rho\in\operatorname{Irr}C(\gamma,u)}
    \rho\otimes K^A(\operatorname{Sp}_u)_{(\gamma,q,\rho)}.
\]
可以证明（细节见 [Kazhdan and Lusztig 1987; Chriss and Ginzburg 1997]），每个 $H(G(F),I)$-模 $K^A(\operatorname{Sp}_u)_{(\gamma,q,\rho)}$ 都有唯一的不可约商；这样就用 (10.4-6) 左端出现的三元组参数化了不可约模。因此同一个集合也与 (10.4-6) 的右端双射。温和局部 Langlands 对应 (10.4-6)，也称 Deligne--Langlands 猜想，就是这样证明的。

\subsection{容许 \texorpdfstring{$(\hat{\fg}_{\kappa_c},I)$}{(\hat g_kc,I)} Harish--Chandra 模的范畴}

\begin{keypoints}
    \item \textbf{核心思想：} 在范畴情形，问题变成：对哪些 $\chi$，范畴 $\hat{\fg}_{\kappa_c}\text{-}\mathrm{mod}_\chi$ 含非零的 $(\hat{\fg}_{\kappa_c},I)$ Harish--Chandra 模？
    \item \textbf{数学工具：} 正则奇点 opers 空间 $\Op^{\mathrm{RS}}_{{}^L G}(D)_{\varpi}$；整权集 $P$；Verma 模与中心特征标（定理 9.5.3）；留数映射。
    \item \textbf{关键结论：} 引理 10.4.1 给出判据：$\chi$ 落在 (10.4-9) 所示的一族正则奇点 opers 中；引理 10.4.2 说明这些 $\chi$ 底层的局部系恰是“温和分歧”的（正则奇点 + 幂幺单值性）。
    \item \textbf{记号提醒：} $I^0=[I,I]$，$I/I^0=H$；$\varpi:{}^L\fh\to{}^L\fh/W$；$\Loc^{\mathrm{RS,tame}}_{{}^L G}\cong\mathcal{N}/{}^L G$。
\end{keypoints}

现在我们希望在圈群范畴 Langlands 对应的框架中，为上述结果寻找范畴类比物。

如 10.2.2 节所解释，在范畴情形，$G(F)$ 的表示被带 $G((t))$ 作用的范畴 $\hat{\fg}_{\kappa_c}\text{-}\mathrm{mod}_\chi$ 替代，而 $I$-不变向量空间被 $\hat{\fg}_{\kappa_c}\text{-}\mathrm{mod}_\chi$ 中 $(\hat{\fg}_{\kappa_c},I)$ Harish--Chandra 模的子范畴替代。于是“哪些 $G(F)$ 的表示容许非零 $I$-不变向量”这一问题，类比为下述问题：对哪些 $\chi$，范畴 $\hat{\fg}_{\kappa_c}\text{-}\mathrm{mod}_\chi$ 含非零的 $(\hat{\fg}_{\kappa_c},I)$ Harish--Chandra 模？答案由下述引理给出。

回忆 9.1.2 节引入了正则奇点且留数为 $\varpi\in{}^L\fh/W=\fh^*/W$ 的 $D$ 上 ${}^L G$-opers 空间 $\Op^{\mathrm{RS}}_{{}^L G}(D)_{\varpi}$，其中 $W$ 是 ${}^L G$ 的外尔群。给定 $\mu\in\fh^*$，用 $\varpi(\mu)$ 记 $\mu$ 在 $\fh^*/W$ 上的投影。最后，设 $P$ 是 $\fg$ 的整权（不必优）的集合，看作 $\fh^*$ 的子集。

\begin{lemma}\label{lem:u16-tame-crit}
    范畴 $\hat{\fg}_{\kappa_c}\text{-}\mathrm{mod}_\chi$ 含非零的 $(\hat{\fg}_{\kappa_c},I)$ Harish--Chandra 模，当且仅当
    \begin{equation}
        \chi\in\bigcup_{\nu\in P/W}\Op^{\mathrm{RS}}_{{}^L G}(D)_{\varpi(\nu)}.
        \tag{10.4-9}
    \end{equation}
\end{lemma}

\begin{proof}
    设 $M$ 是 $\hat{\fg}_{\kappa_c}\text{-}\mathrm{mod}_\chi$ 中的非零 $(\hat{\fg}_{\kappa_c},I)$ Harish--Chandra 模。我们按引理 10.3.2 证明中的同样方式证明，$M$ 含一个被 $I^0=[I,I]$ 零化、且使 $I/I^0=H$ 按特征标 $\nu\in P$ 作用的向量 $v$。因此存在从 Verma 模 $\mathbb{M}_\nu$ 到 $M$ 的非零同态，把 $\mathbb{M}_\nu$ 的最高权向量送到 $v$。根据定理 9.5.3，中心 $Z(\hat{\fg})\simeq\operatorname{Fun}\Op_{{}^L G}(D^\times)$ 在 $\mathbb{M}_\nu$ 上的作用经过 $\operatorname{Fun}\Op^{\mathrm{RS}}_{{}^L G}(D)_{\varpi(-\nu-\rho)}$ 分解。因此 $\chi\in\Op^{\mathrm{RS}}_{{}^L G}(D)_{\varpi(-\nu-\rho)}$。

    另一方面，设 $\chi$ 属于某个 $\nu\in P$ 的 $\Op^{\mathrm{RS}}_{{}^L G}(D)_{\varpi(\nu)}$。设 $\mathbb{M}_\nu(\chi)$ 是 Verma 模 $\mathbb{M}_\nu$ 模去对应于 $\chi$ 的中心特征标所得的商。由定理 9.5.3，$\mathbb{M}_\nu(\chi)$ 是 $\hat{\fg}_{\kappa_c}\text{-}\mathrm{mod}_\chi$ 的非零对象。另一方面，显然 $\mathbb{M}_\nu$、从而 $\mathbb{M}_\nu(\chi)$，是 $I$-等变的 $\hat{\fg}_{\kappa_c}$-模。
\end{proof}

% [Tutor Note: p.343 原文内部不一致，译文照原书保留。引理 10.4.1 的陈述 (10.4-9) 写 $\varpi(\nu)$（$\nu\in P/W$），但证明第一段得到的是 $\varpi(-\nu-\rho)$，第二段又用 $\varpi(\nu)$。两个参数相差一个“平移”，不可能同时正确：按 Verma 模的中心特征标惯例，正确形式应以 $\varpi(-\nu-\rho)$（即 $-\nu-\rho$ 的投影）为准，或陈述中的 $\nu$ 应理解为 $\nu+\rho$ 的平移。此点已登记报告，请引用时注意。]
于是，使相应范畴 $\hat{\fg}_{\kappa_c}\text{-}\mathrm{mod}_\chi$ 含非平凡 $I$-等变对象的那些 opers $\chi$，正是子概形 (10.4-9) 的点。\transnote{此处 (10.4-9) 的平移与证明不一致：陈述与证明第二段用 $\varpi(\nu)$（$\nu\in P$ 或 $P/W$），证明第一段却得到 $\varpi(-\nu-\rho)$。按定理 9.5.3（Verma 模 $\mathbb M_\nu$ 的中心特征标对应 $\Op^{\mathrm{RS}}_{{}^L G}(D)_{\varpi(-\nu-\rho)}$），一致形式应为 $\varpi(-\nu-\rho)$；陈述中的 $\varpi(\nu)$ 系原书笔误或另一平移约定，本译本照原文保留、未改公式。引用时请以定理 9.5.3 为准；本译注对应遗留问题 A53。[来源：本书 p.343 与 p.287–288 定理 9.5.3；[Frenkel and Gaitsgory 2006c, Corollary 13.3.2]（本书 p.288 引用）。]}下一个问题是：它们对应的 ${}^L G$-局部系是什么？

设 $\Loc^{\mathrm{RS,tame}}_{{}^L G}\subset\Loc_{{}^L G}(D^\times)$ 是 $D^\times$ 上具正则奇点与幂幺单值性的 ${}^L G$-局部系的轨迹。这样的局部系在同构意义下由其单值性的共轭类决定（例如见 [Babbitt and Varadarajan 1983] 第 8 节）。因此 $\Loc^{\mathrm{RS,tame}}_{{}^L G}$ 是代数栈，同构于 $\mathcal{N}/{}^L G$。下述结果的证明与引理 10.3.1 的证明类似。

\begin{lemma}\label{lem:u16-tame-pre}
    若 oper $\chi\in\Op_{{}^L G}(D^\times)$ 底层的局部系属于 $\Loc^{\mathrm{RS,tame}}_{{}^L G}$，则 $\chi$ 属于 $\Op_{{}^L G}(D^\times)$ 的子集 (10.4-9)。
\end{lemma}

换言之，子概形 (10.4-9) 恰是 $\Loc^{\mathrm{RS,tame}}_{{}^L G}\subset\Loc_{{}^L G}(D^\times)$ 在映射 $\alpha:\Op_{{}^L G}(D^\times)\to\Loc_{{}^L G}(D^\times)$ 下的（集合论意义下的）原像。

这并不令人意外。事实上，类比经典 Langlands 对应，我们期待含非平凡 $I$-等变对象的范畴 $\hat{\fg}_{\kappa_c}\text{-}\mathrm{mod}_\chi$ 对应那些作为温和分歧同态 $W'_F\to{}^L G$ 的几何对应物的 Langlands 参数。这些参数最显然的候选，正是 $D^\times$ 上具正则奇点与幂幺单值性的 ${}^L G$-局部系。因此我们把这样的局部系称为温和分歧的（这反映在上面的记号中）。

% [Tutor Note: p.344。这里的逻辑链值得记住：$I$-等变对象存在 ⟺ $\chi$ 在 (10.4-9) 中（引理 10.4.1）；而 (10.4-9) 恰是“正则奇点 + 幂幺单值性”的原像（引理 10.4.2）。于是“温和分歧”一词在几何一侧有了精确定义：底层局部系的单值性幂幺。注意单值性幂幺等价于留数幂零，故 $\Loc^{\mathrm{RS,tame}}\cong\mathcal{N}/{}^L G$。]
总结一下：设 $\sigma$ 是 $D^\times$ 上的温和分歧 ${}^L G$-局部系，$\chi$ 是与 $\sigma$ 规范等价的 ${}^L G$-oper。则按引理 10.4.1，$\chi$ 属于子概形 (10.4-9)，且相应的范畴 $\hat{\fg}_{\kappa_c}\text{-}\mathrm{mod}_\chi$ 含非零的 $I$-等变对象。设 $\hat{\fg}_{\kappa_c}\text{-}\mathrm{mod}_\chi^I$ 是相应的 $I$-等变（等价地，$(\hat{\fg}_{\kappa_c},I)$ Harish--Chandra）模范畴。这就是我们对“温和分歧同态 $W'_F\to{}^L G$ 对应的不可约 $G(F)$ 表示中 $I$-不变向量空间”的范畴化候选。

注意根据 猜想 10.1.2，对一切作为 ${}^L G$-局部系彼此规范等价的 $\chi$，范畴 $\hat{\fg}_{\kappa_c}\text{-}\mathrm{mod}_\chi$（相应地，$\hat{\fg}_{\kappa_c}\text{-}\mathrm{mod}_\chi^I$）应当彼此等价。

在下一小节，我们按 [Frenkel and Gaitsgory 2006c] 用 Springer 纤维上凝聚层的范畴，对 $\chi\in\Op^{\mathrm{RS}}_{{}^L G}(D)_{\varpi(-\rho)}$ 的范畴 $\hat{\fg}_{\kappa_c}\text{-}\mathrm{mod}_\chi^I$ 给出猜想性描述。这个描述特别蕴含：至少对这些范畴的导出范畴，对应于同一局部系的 opers 给出相互等价的范畴。对 $\chi\in\Op^{\mathrm{RS}}_{{}^L G}(D)_{\varpi(\nu)}$（其他 $\nu\in P$）也有类似的猜想，读者可仿照我们对 $\nu=-\rho$ 情形的讨论自行写出。

\subsection{\texorpdfstring{$(\hat{\fg}_{\kappa_c},I)$}{(g_kc,I)} Harish--Chandra 模范畴的猜想描述}

\begin{keypoints}
    \item \textbf{核心思想：} 在温和分歧情形，阿贝尔范畴不够用，必须改用导出范畴 $D^b(\hat{\fg}_{\kappa_c}\text{-}\mathrm{mod}_\chi)^I$；猜想它等价于某个（加厚、DG 化）Springer 纤维上拟凝聚层的导出范畴。
    \item \textbf{数学工具：} $I$-monodromic 模（$I^0$-等变替代）；非平坦的 Springer 分解导致的高阶 Tor；DG 纤维与 DG 概形；Grothendieck alteration 与加厚 Springer 分解；Miura opers 空间 $\widetilde{\mathrm{MOp}}{}^0_{{}^L G}$。
    \item \textbf{关键结论：} 主猜想 10.4.4：(10.4-13) 式的范畴等价，且它\textbf{不保 $t$-结构}；对单个 oper 的特殊化给出 (10.4-14)。
    \item \textbf{记号提醒：} $I^0=[I,I]$；$D^b(\hat{\fg}_{\kappa_c}\text{-}\mathrm{mod}_{\mathrm{nilp}})^{I^0}$ 是“族版本”；$\widetilde{\mathrm{MOp}}{}^0_{{}^L G}$ 是幂零 Miura opers 的族。
\end{keypoints}

% [Tutor Note: p.344--345。本小节是全书技术性最强的地方之一：三层“加厚”依次出现——(1) 从 $I$ 到 $I^0$（$I$-单值模），(2) 从普通纤维到 DG 纤维（因 Springer 分解非平坦），(3) 从约化纤维到非约化（加厚）Springer 纤维。最终用“族版本”避开 DG 概形，得到干净的 猜想 10.4.4。]
让我们考虑子概形 (10.4-9) 的一个连通分支，即 $\Op^{\mathrm{RS}}_{{}^L G}(D)_{\varpi(-\rho)}$。回忆 9.2.1 节与命题 9.2.1：这个空间同构于 $D$ 上幂零 opers 的空间 $\Op^{\mathrm{nilp}}_{{}^L G}$（对应权 $\lambda=0$）。如 9.2.3 节所解释，它带有留数映射
\[
    \Res_{\mathcal{F}}:\Op^{\mathrm{nilp}}_{{}^L G}\to{}^L\fn_{\mathcal{F}_{{}^L B,0}},
    \qquad
    \Res:\Op^{\mathrm{nilp}}_{{}^L G}\to{}^L\fn/{}^L B=\widetilde{\mathcal{N}}/{}^L G.
\]
对任意 $\chi\in\Op^{\mathrm{nilp}}_{{}^L G}$，相应联络的 ${}^L G$-规范等价类是 $D^\times$ 上的温和分歧 ${}^L G$-局部系。此外，其单值性共轭类等于 $\exp(2\pi i\Res(\chi))$。

我们希望描述中心特征标为 $\chi\in\Op^{\mathrm{nilp}}_{{}^L G}$ 的 $(\hat{\fg}_{\kappa_c},I)$ Harish--Chandra 模范畴 $\hat{\fg}_{\kappa_c}\text{-}\mathrm{mod}_\chi^I$。然而这里遇到与无分歧情形相比的第一个重大困难。虽然在无分歧情形我们用的是阿贝尔范畴 $\hat{\fg}_{\kappa_c}\text{-}\mathrm{mod}_\chi^{G[[t]]}$，在温和分歧情形这似乎不可能。因此从现在起我们将使用适当的导出范畴 $D^b(\hat{\fg}_{\kappa_c}\text{-}\mathrm{mod}_\chi)^I$。按定义，它是 $\hat{\fg}_{\kappa_c}\text{-}\mathrm{mod}_\chi$ 的有界导出范畴 $D^b(\hat{\fg}_{\kappa_c}\text{-}\mathrm{mod}_\chi)$ 的满子范畴，其对象是上同调落在 $\hat{\fg}_{\kappa_c}\text{-}\mathrm{mod}_\chi^I$ 中的复形。

粗略地说，[Frenkel and Gaitsgory 2006c] 的猜想是：范畴 $D^b(\hat{\fg}_{\kappa_c}\text{-}\mathrm{mod}_\chi)^I$ 等价于拟凝聚层范畴 $\mathrm{QCoh}(\operatorname{Sp}_{\Res_{\mathcal{F}}(\chi)})$ 的导出范畴 $D^b(\mathrm{QCoh}(\operatorname{Sp}_{\Res_{\mathcal{F}}(\chi)}))$，其中 $\operatorname{Sp}_{\Res_{\mathcal{F}}(\chi)}$ 是 $\Res_{\mathcal{F}}(\chi)$ 的 Springer 纤维。然而需要对这个陈述作一些调整。这些调整是为了得到下文的“漂亮”陈述，即 猜想 10.4.4。下面解释这些调整是什么、以及为什么需要它们。\transnote{三个调整各有原因：$I$-monodromic 模之间的扩张未必 $I$-等变，故先放宽到 $I^0$-等变；Springer 分解 $\widetilde{\mathcal N}\to\mathcal N$ 非平坦，普通张量积只给出约化纤维，故须取导出张量积与 DG 纤维；最后把约化纤维换成“加厚”纤维，才能使每点的纤维是 $\Spec h_0$ 而不是单点。三者依次放宽后，陈述才能用族版本（避开 DG 概形）写成干净的 猜想 10.4.4。[来源：本书 p.344–346；[Frenkel and Gaitsgory 2006c]（本书 p.344 引用）。]}

第一个调整是：需要考虑比 $D^b(\hat{\fg}_{\kappa_c}\text{-}\mathrm{mod}_\chi)^I$ 稍大的表示范畴。即，我们希望包含 $I$-等变 $\hat{\fg}_{\kappa_c}$-模的扩张，这些扩张不必是 $I$-等变的，而只是 $I^0$-等变的，其中 $I^0=[I,I]$。为更精确地解释，取 Cartan 子群 $H\subset B\subset I$ 及相应的李子代数 $\fh\subset\fb\subset\Lie I$。于是有同构 $I=H\ltimes I^0$。一个 $I$-等变 $\hat{\fg}_{\kappa_c}$-模，等同于一个模，其上 $\fh$ 按整权给出的特征值对角地作用，且李代数 $\Lie I^0$ 局部幂零地作用。然而这样的模之间可能存在扩张，在扩张上 $\fh$ 的作用不再半单。这样的模称为 $I$-monodromic。更精确地说，$I$-monodromic 的 $\hat{\fg}_{\kappa_c}$-模是指具有递增滤过、其相邻商为 $I$-等变的模。把这样的模纳入我们的范畴是自然的。但容易证明：对任意 $\chi\in\Op^{\mathrm{nilp}}_{{}^L G}$，$\hat{\fg}_{\kappa_c}\text{-}\mathrm{mod}_\chi$ 中的 $I$-monodromic 对象等同于 $I^0$-等变对象（见 [Frenkel and Gaitsgory 2006c]）。因此我们将用 $I^0$-等变模代替 $I$-monodromic 模。用 $D^b(\hat{\fg}_{\kappa_c}\text{-}\mathrm{mod}_\chi)^{I^0}$ 记 $D^b(\hat{\fg}_{\kappa_c}\text{-}\mathrm{mod}_\chi)$ 中上同调落在 $\hat{\fg}_{\kappa_c}\text{-}\mathrm{mod}_\chi^{I^0}$ 里的对象组成的满子范畴。

第二个调整涉及 Springer 分解 $\widetilde{\mathcal{N}}\to\mathcal{N}$ 的非平坦性。按定义，Springer 纤维 $\operatorname{Sp}_u$ 是纤维积 $\widetilde{\mathcal{N}}\times_{\mathcal{N}}\mathrm{pt}$，其中 $\mathrm{pt}$ 是点 $u\in\mathcal{N}$。这意味着 $\operatorname{Sp}_u$ 的结构层由
\begin{equation}
    \cO_{\operatorname{Sp}_u}=\cO_{\widetilde{\mathcal{N}}}\otimes_{\cO_{\mathcal{N}}}\C
    \tag{10.4-10}
\end{equation}
给出。然而由于态射 $\widetilde{\mathcal{N}}\to\mathcal{N}$ 不是平坦的，这个张量积函子不是左正合的，并且存在非平凡的导出张量积（Tor）。我们（猜想中的）等价不会是正合函子：它把范畴 $\hat{\fg}_{\kappa_c}\text{-}\mathrm{mod}_\chi^{I^0}$ 的一般对象送到拟凝聚层范畴的复形，更精确地说，送到相应导出范畴的对象。因此我们被迫使用导出范畴，从而必须计入高阶导出张量积。

% [Tutor Note: p.345--346 易错处。(10.4-10) 中 $\cO_{\widetilde{\mathcal{N}}}\otimes_{\cO_{\mathcal{N}}}\C$ 若按普通张量积理解只给出约化纤维；要得到正确的（非约化）纤维必须取导出张量积。这是“非平坦”造成的：$\widetilde{\mathcal{N}}\to\mathcal{N}$ 在奇点处纤维维数跳变，Tor 非零。作者随后用“DG 纤维”与“加厚”两种语言描述同一现象，读者需对照 (10.4-11) 与 (10.4-13)。]
为更好地理解这种非正合性的后果，考虑如下模型例子。设我们已建立导出范畴 $D^b(\mathrm{QCoh}(\widetilde{\mathcal{N}}))$ 与另一个导出范畴 $D^b(\mathcal{C})$ 之间的等价。特别地，这意味着两个范畴都带代数 $\operatorname{Fun}\mathcal{N}$ 的作用（回忆 $\mathcal{N}$ 是仿射代数簇）。设 $\operatorname{Fun}\mathcal{N}$ 在 $D^b(\mathcal{C})$ 上的作用来自它在阿贝尔范畴 $\mathcal{C}$ 上的作用。于是 $\mathcal{C}$ 纤维化于 $\mathcal{N}$ 之上，设 $\mathcal{C}_u$ 是对应于 $u\in\mathcal{N}$ 的纤维范畴。这是 $\mathcal{C}$ 的满子范畴，其对象是 $\mathcal{C}$ 中在 $\operatorname{Fun}\mathcal{N}$ 里 $u$ 的理想作用下为零的对象。\footnote{$\mathcal{C}$ 与 $\mathcal{C}_u$ 的关系类似于 $\hat{\fg}_{\kappa_c}\text{-}\mathrm{mod}$ 与 $\hat{\fg}_{\kappa_c}\text{-}\mathrm{mod}_\chi$ 的关系，其中 $\chi\in\Op_{{}^L G}(D^\times)$。} 范畴 $D^b(\mathcal{C}_u)$ 等价于什么？

容易想说它等价于 $D^b(\mathrm{QCoh}(\operatorname{Sp}_u))$。然而这并不能由 $D^b(\mathrm{QCoh}(\mathcal{N}))$ 与 $D^b(\mathcal{C})$ 的等价推出，因为张量积 (10.4-10) 有非平凡的高阶导出函子。设范畴 $\mathcal{C}$ 在 $\mathcal{N}$ 上是平坦的，即 $\mathcal{C}$ 的所有投射对象作为 $\operatorname{Fun}\mathcal{N}$-模都是平坦的。则 $D^b(\mathcal{C}_u)$ 等价于范畴 $D^b(\mathrm{QCoh}(\operatorname{Sp}_u^{\mathrm{DG}}))$，其中 $\operatorname{Sp}_u^{\mathrm{DG}}$ 是 $\widetilde{\mathcal{N}}\to\mathcal{N}$ 在 $u$ 处的“DG 纤维”。按定义，$\operatorname{Sp}_u^{\mathrm{DG}}$ 上的拟凝聚层是 DG 代数
\begin{equation}
    \cO_{\operatorname{Sp}_u^{\mathrm{DG}}}
    =\cO_{\widetilde{\mathcal{N}}}\overset{L}{\otimes}_{\cO_{\mathcal{N}}}\C_u
    \tag{10.4-11}
\end{equation}
上的 DG 模，其中现在取张量积的完全导出函子。于是范畴 $D^b(\mathrm{QCoh}(\operatorname{Sp}_u^{\mathrm{DG}}))$ 可以看作“DG 概形” $\operatorname{Sp}_u^{\mathrm{DG}}$ 上拟凝聚层的导出范畴（DG 概形的精确定义见 [Ciocan-Fontanine and Kapranov 2001]）。

最后，第三个调整是应考虑非约化的 Springer 纤维。这意味着不用 Springer 分解 $\widetilde{\mathcal{N}}$，而用“加厚”的 Springer 分解
\[
    {}^{L}\widetilde{\widetilde{\mathcal{N}}}
    ={}^{L}\widetilde{\fg}\times_{{}^L\fg}\mathcal{N},
\]
其中 ${}^L\widetilde{\fg}$ 是所谓 Grothendieck alteration\transnote{Grothendieck alteration 指 ${}^L\widetilde{\fg}=\{(x,\mathfrak b')\mid x\in\mathfrak b'\}\to{}^L\fg$：它在正则半单元上的纤维是 $|W|$ 个点，因此是 de Jong 意义下“紧合、纤维有限”的 alteration，而不是双有理态射。把它限制到幂零锥得到的加厚版本 ${}^L\widetilde{\widetilde{\mathcal N}}$ 在正则幂零元上的纤维才是非约化的 $\Spec h_0$。本书按术语表保留此英文原名。[来源：本书 p.346–347；[Chriss and Ginzburg 1997]（本书 p.345 引用）。]}
\[
    {}^{L}\widetilde{\fg}=\{x\in{}^L\fg,\ \mathfrak{b}'\in{}^L G/{}^L B\mid x\in\mathfrak{b}'\},
\]
我们在 7.1.1 节已遇到过。代数簇 ${}^L\widetilde{\widetilde{\mathcal{N}}}$ 非约化，其底层的约化簇是 Springer 分解 $\widetilde{\mathcal{N}}$。例如，$\widetilde{\mathcal{N}}$ 在 $\mathcal{N}$ 的正则元上的纤维是单点，但 ${}^L\widetilde{\widetilde{\mathcal{N}}}$ 的相应纤维是 Artin 代数
\[
    h_0=\operatorname{Fun}{}^L\fh\big/\bigl(\operatorname{Fun}{}^L\fh\bigr)^W_+
\]
的谱。这里 $(\operatorname{Fun}{}^L\fh)^W_+$ 是 $\operatorname{Fun}{}^L\fh$ 中由 $W$-不变子代数的增广理想生成的理想。因此 $\Spec h_0$ 是 $\varpi:{}^L\fh\to{}^L\fh/W$ 在 $0$ 处的概形论纤维。\transnote{$h_0=\operatorname{Fun}{}^L\fh/(\operatorname{Fun}{}^L\fh)^W_+$ 是 $W$ 的余不变量代数：作为 $W$-模它同构于正则表示，故 $\dim h_0=|W|$（Chevalley 定理）；而 $\Spec h_0$ 恰是 $\varpi$ 在 $0$ 处的概形论纤维。这解释了加厚纤维的“长度”为什么恰好是 $|W|$，也预示了 10.4.6 节 $\fsl_2$ 情形的双点 $h_0=\C[\epsilon]/(\epsilon^2)$。[来源：本书 p.347；[Frenkel and Gaitsgory 2006c]（本书 p.345–347 引用）。]}结果表明，为描述范畴 $D^b(\hat{\fg}_{\kappa_c}\text{-}\mathrm{mod}_\chi)^{I^0}$，我们需要使用“加厚”的 Springer 分解。\footnote{这一点在 10.4.6 节对 $\fg=\fsl_2$ 作了解释。}

% [Tutor Note: p.347。加厚分解的作用：把“$\chi$ 的留数为正则幂零元时纤维是一个点”与“应该得到 $\Spec h_0$”这两个事实协调起来。$h_0$ 的维数是 $|W|$，这正是 外尔群大小的来源，也是 Hecke 代数中心 $\Rep{}^L G$ 的“外尔群对称性”的几何体现。]
总结一下：为构造所寻求的范畴等价，我们不取单个 Springer 纤维，而取整个 Springer 分解，并进一步把它换成上面定义的“加厚”Springer 分解 ${}^L\widetilde{\widetilde{\mathcal{N}}}$。在这个版本里，我们能表述避开 DG 概形的等价。

这意味着：不应考虑单个幂零 oper $\chi$ 对应的范畴 $\hat{\fg}_{\kappa_c}\text{-}\mathrm{mod}_\chi$，而应考虑“万有”范畴 $\hat{\fg}_{\kappa_c}\text{-}\mathrm{mod}_{\mathrm{nilp}}$，它是所有这些范畴的“族版本”。按定义，范畴 $\hat{\fg}_{\kappa_c}\text{-}\mathrm{mod}_{\mathrm{nilp}}$ 是 $\hat{\fg}_{\kappa_c}\text{-}\mathrm{mod}$ 的满子范畴，其对象的性质是：$Z(\hat{\fg})=\operatorname{Fun}\Op_{{}^L G}(D)$ 在它们上的作用经过商 $\operatorname{Fun}\Op_{{}^L G}(D)\to\operatorname{Fun}\Op^{\mathrm{nilp}}_{{}^L G}$ 分解。因此范畴 $\hat{\fg}_{\kappa_c}\text{-}\mathrm{mod}_{\mathrm{nilp}}$ 类似于上面考虑过的范畴 $\hat{\fg}_{\kappa_c}\text{-}\mathrm{mod}_{\mathrm{reg}}$。前者纤维化于 $\Op^{\mathrm{nilp}}_{{}^L G}$ 之上，后者纤维化于 $\Op_{{}^L G}(D)$ 之上。各个范畴 $\hat{\fg}_{\kappa_c}\text{-}\mathrm{mod}_\chi$ 现在实现为这些范畴在特定 opers $\chi$ 上的纤维。

我们天真的想法是：对每个 $\chi\in\Op^{\mathrm{nilp}}_{{}^L G}$，范畴 $D^b(\hat{\fg}_{\kappa_c}\text{-}\mathrm{mod}_\chi)^{I^0}$ 等价于 $\mathrm{QCoh}(\operatorname{Sp}_{\Res_{\mathcal{F}}(\chi)})$。我们现在想表述这种等价的一个“族版本”。为此构造纤维积
\[
    {}^{L}\widetilde{\widetilde{\mathrm{n}}}
    ={}^{L}\widetilde{\fg}\times_{{}^L\fg}{}^L\mathfrak{n}.
\]
结果表明这个纤维积不受单个 Springer 纤维的问题困扰，如下引理所示：

\begin{lemma}[{[Frenkel and Gaitsgory 2006c], 引理 6.4}]\label{lem:u16-flat}
    导出张量积
    \[
        \operatorname{Fun}{}^L\widetilde{\fg}
        \overset{L}{\otimes}_{\operatorname{Fun}{}^L\fg}
        \operatorname{Fun}{}^L\mathfrak{n}
    \]
    集中在上同调维数 $0$。
\end{lemma}

% [Tutor Note: p.347--348。引理 10.4.3 是全书少见的“技术性引理直接决定主猜想形式”的例子：因为 $\operatorname{Fun}{}^L\widetilde{\fg}\otimes^L_{\operatorname{Fun}{}^L\fg}\operatorname{Fun}{}^L\mathfrak{n}$ 集中在一处，族版本的等价才能不用 DG 概形而写成普通导出范畴。注意引理原文的记号是 $\operatorname{Fun}{}^L\widetilde{\fg}$（不是 $\cO$），这是“函数环”版本的 formulation。]
代数簇 ${}^L\widetilde{\widetilde{\mathrm{n}}}$ 可以看作以 ${}^L\mathfrak{n}\subset{}^L\fg$ 为参数族的（非约化）Springer 纤维的族。重要的是指出它奇异、可约且非约化。例如若 $\fg=\fsl_2$，它有两个分支，其一是 $\PP^1$（$0$ 处的 Springer 纤维），其二是加倍的仿射线（即 $\Spec\C[x,y]/(y^2)$），见 10.4.6 节。

我们指出相应的约化概形是 9.3.2 节引入的 ${}^L\widetilde{\mathrm{n}}$：
\begin{equation}
    {}^{L}\widetilde{\mathrm{n}}=\widetilde{\mathcal{N}}\times_{\mathcal{N}}{}^L\mathfrak{n}.
    \tag{10.4-12}
\end{equation}
然而对应于 (10.4-12) 的导出张量积并不集中在上同调维数 $0$，这就是我们宁用 ${}^L\widetilde{\widetilde{\mathrm{n}}}$ 而不用 ${}^L\widetilde{\mathrm{n}}$ 的原因。

现在令
\[
    \widetilde{\mathrm{MOp}}{}^0_{{}^L G}
    =\Op^{\mathrm{nilp}}_{{}^L G}
    \times_{{}^L\mathfrak{n}/{}^L B}
    {}^{L}\widetilde{\widetilde{\mathrm{n}}}/{}^L B,
\]
其中用了留数态射 $\Res:\Op^{\mathrm{nilp}}_{{}^L G}\to{}^L\mathfrak{n}/{}^L B$。于是可以非正式地把 $\widetilde{\mathrm{MOp}}{}^0_{{}^L G}$ 想成 $\Op^{\mathrm{nilp}}_{{}^L G}$ 上的族，它在 $\chi\in\Op^{\mathrm{nilp}}_{{}^L G}$ 上的纤维是 $\Res(\chi)$ 的 Springer 纤维。

正如记号所示，$\widetilde{\mathrm{MOp}}{}^0_{{}^L G}$ 与底层 opers 为幂零的 Miura opers 空间密切相关（见下文 10.4.5 节）。更精确地说，$\widetilde{\mathrm{MOp}}{}^0_{{}^L G}$ 的约化概形是 9.3.2 节引入的权 $0$ 的幂零 Miura ${}^L G$-opers 的概形 $\mathrm{MOp}^0_{{}^L G}$：
\[
    \mathrm{MOp}^0_{{}^L G}
    \simeq\Op^{\mathrm{nilp}}_{{}^L G}\times_{{}^L\mathfrak{n}/{}^L B}{}^L\widetilde{\mathrm{n}}/{}^L B
\]
（见公式 (9.3-9) 与 (10.4-12)）。

我们还引入范畴 $\hat{\fg}_{\kappa_c}\text{-}\mathrm{mod}_{\mathrm{nilp}}^{I^0}$，即 $\hat{\fg}_{\kappa_c}\text{-}\mathrm{mod}_{\mathrm{nilp}}$ 中由 $I^0$-等变对象组成的满子范畴，以及相应的导出范畴
\[
    D^b\bigl(\hat{\fg}_{\kappa_c}\text{-}\mathrm{mod}_{\mathrm{nilp}}\bigr)^{I^0}.
\]
现在可以陈述 [Frenkel and Gaitsgory 2006c] 的主猜想了：\footnote{校样时补注：猜想 10.4.4 已由 E. Frenkel 与 D. Gaitsgory 在 [Frenkel and Gaitsgory 2007a] 中证明。}

\begin{conjecture}\label{conj:u16-main}
    存在范畴等价
    \begin{equation}
        D^b\bigl(\hat{\fg}_{\kappa_c}\text{-}\mathrm{mod}_{\mathrm{nilp}}\bigr)^{I^0}
        \simeq D^b\bigl(\mathrm{QCoh}(\widetilde{\mathrm{MOp}}{}^0_{{}^L G})\bigr),
        \tag{10.4-13}
    \end{equation}
    它与代数 $\operatorname{Fun}\Op^{\mathrm{nilp}}_{{}^L G}$ 在两个范畴上的作用相容。
\end{conjecture}

注意 $\operatorname{Fun}\Op^{\mathrm{nilp}}_{{}^L G}$ 在第一个范畴上的作用来自中心 $Z(\hat{\fg})$ 的作用，在第二个范畴上的作用则来自 $\widetilde{\mathrm{MOp}}{}^0_{{}^L G}$ 是 $\Op^{\mathrm{nilp}}_{{}^L G}$ 上的概形这一事实。

% [Tutor Note: p.348 主猜想。猜想 10.4.4 是本节的核心陈述，可视为“温和分歧局部 Langlands 对应的范畴版本”。它把 $\hat{\fg}_{\kappa_c}$ 临界水平模范畴的导出范畴与一个几何对象（加厚 Miura opers 族）上的拟凝聚层导出范畴等同；“与 $\operatorname{Fun}\Op^{\mathrm{nilp}}$ 的作用相容”这一条件保证等价沿 $\Op^{\mathrm{nilp}}$ 的纤维化是逐点相容的。]
猜想性等价 (10.4-13) 可以看作温和分歧情形局部 Langlands 对应的范畴类比物，下一小节将讨论这一点。\transnote{猜想 10.4.4 是 [Frenkel and Gaitsgory 2006c] 的主猜想；据校样补注（见前一脚注），它已由作者在 [Frenkel and Gaitsgory 2007a] 中证明。注意等价 (10.4-13) 不保 $t$-结构，故“阿贝尔范畴层面”的直觉不能直接搬运；这正是全节坚持导出范畴、并在单点特殊化 (10.4-14) 中出现 DG Springer 纤维的原因。陈述中“与 $\operatorname{Fun}\Op^{\mathrm{nilp}}$ 的作用相容”这一条同样是猜想的一部分。[来源：本书 p.348 及同页脚注；[Frenkel and Gaitsgory 2007a]（本书 p.348 脚注引用）。]}

等价 (10.4-13) 的一个重要特征是它不保持两个范畴上的 $t$-结构。换言之，阿贝尔范畴 $\hat{\fg}_{\kappa_c}\text{-}\mathrm{mod}_{\mathrm{nilp}}^{I^0}$ 的对象在此等价下一般映为 $D^b(\mathrm{QCoh}(\widetilde{\mathrm{MOp}}{}^0_{{}^L G}))$ 中的复形，反之亦然。我们将在 10.4.6 节 $\fsl_2$ 的情形看到这样的例子。

对应于带整优权 $\lambda\in P^+$ 的幂零 opers 空间 $\Op^{\mathrm{nilp},\lambda}_{{}^L G}$ 的范畴，也有类似的猜想。

在下一小节，我们将讨论 猜想 10.4.4 与经典温和分歧 Langlands 对应之间的联系。然后给出这个猜想的一些证据，并详细考察 $\fg=\fsl_2$ 的例子。

\subsection{经典情形与几何情形的联系}

\begin{keypoints}
    \item \textbf{核心思想：} 把 (10.4-13) 特殊化到单个 oper，得到 (10.4-14)；再用 Grothendieck 群把它与经典 $K$-理论、从而与仿射 Hecke 代数的模联系起来。
    \item \textbf{数学工具：} DG Springer 纤维的 Grothendieck 群等于 $K(\operatorname{Sp}_u)$；$K^A(\operatorname{Sp}_u)_{(\gamma,q)}$ 与 $\operatorname{Hom}_{C(\gamma,u)}(\rho,K(\operatorname{Sp}_u))$；$\Rep{}^L G$ 的范畴化与纤维函子；$^L G$-torsor。
    \item \textbf{关键结论：} 经典参数 $\gamma$ 与 $\rho$ 在几何一侧被“局部系的对称群 $\operatorname{Aut}(\sigma)$ 的作用”吸收；$q$ 的取值在几何情形取 $1$。
    \item \textbf{记号提醒：} $\operatorname{Sp}^{\mathrm{DG}}_u$ 是 DG Springer 纤维；$\mathcal{C}_{\mathrm{tame}}$ 是温和分歧 Langlands 范畴；$\operatorname{Aut}(\sigma)$ 是局部系的自同构群。
\end{keypoints}

% [Tutor Note: p.349--350。本小节是“为什么主猜想是对的”的启发式论证：Grothendieck 群层面必须与经典 $K$-理论一致。关键洞察有二：(1) DG/非约化结构在 $K_0$ 中消失，故只剩 $K(\operatorname{Sp}_u)$；(2) 经典参数 $\gamma$（Frobenius）在几何情形没有对应，其角色由局部系的自同构群 $\operatorname{Aut}(\sigma)$ 承担。]
让我们讨论等价 (10.4-13) 与用 Springer 纤维的 $K$-理论实现仿射 Hecke 代数表示之间的联系。如我们已解释，我们希望把 $\chi\in\Op^{\mathrm{nilp}}_{{}^L G}$ 的范畴 $D^b(\hat{\fg}_{\kappa_c}\text{-}\mathrm{mod}_\chi)^{I^0}$ 大致看作不可约表示 $(R,\pi)$ 的 $I$-不变向量空间 $R^{\pi(I)}$ 的范畴化。因此我们期待范畴 $D^b(\hat{\fg}_{\kappa_c}\text{-}\mathrm{mod}_\chi)^{I^0}$ 的 Grothendieck 群与空间 $R^{\pi(I)}$ 有某种联系。

让我们尝试把 猜想 10.4.4 的陈述特殊化到特定的 oper $\chi=(\mathcal{F},\nabla,\mathcal{F}_{{}^L B})\in\Op^{\mathrm{nilp}}_{{}^L G}$。设 $\widetilde{\operatorname{Sp}}^{\mathrm{DG}}_{\Res_{\mathcal{F}}(\chi)}$ 是 $\widetilde{\mathrm{MOp}}{}^0_{{}^L G}$ 在 $\chi$ 上的 DG 纤维。按定义（见 9.2.3 节），$\chi$ 的留数 $\Res_{\mathcal{F}}(\chi)$ 是 ${}^L\mathfrak{n}$ 经 ${}^L B$-torsor $\mathcal{F}_{{}^L B,0}$ 扭变后的一个向量。由此推出 $\widetilde{\operatorname{Sp}}^{\mathrm{DG}}_{\Res_{\mathcal{F}}(\chi)}$ 是 Grothendieck alteration 的 $\mathcal{F}_{{}^L B,0}$-扭变在 $\Res_{\mathcal{F}}(\chi)$ 上的 DG 纤维。

若把 $\mathcal{F}_{{}^L B,0}$ 平凡化，则 $u=\Res_{\mathcal{F}}(\chi)$ 成为 ${}^L\mathfrak{n}$ 的元素。按定义，（非约化的）DG Springer 纤维 $\widetilde{\operatorname{Sp}}^{\mathrm{DG}}_u$ 是 Grothendieck alteration ${}^L\widetilde{\fg}\to{}^L\fg$ 在 $u$ 处的 DG 纤维。换言之，相应的结构层是 DG 代数
\[
    \cO_{\widetilde{\operatorname{Sp}}^{\mathrm{DG}}_u}
    =\cO_{{}^L\widetilde{\fg}}
    \overset{L}{\otimes}_{\cO_{{}^L\fg}}\C_u,
\]
（与公式 (10.4-11) 比较）。

为看清这些 DG 纤维是什么样子，取 $u=0$。则朴素的 Springer 纤维就是旗簇 ${}^L G/{}^L B$（此时它约化），而 $\cO_{\widetilde{\operatorname{Sp}}_0}$ 是 ${}^L G/{}^L B$ 的结构层。但层 $\cO_{\widetilde{\operatorname{Sp}}^{\mathrm{DG}}_0}$ 是 DG 代数层，它拟同构于 ${}^L G/{}^L B$ 上微分形式的复形配以零微分。换言之，$\widetilde{\operatorname{Sp}}^{\mathrm{DG}}_0$ 可以看作一个“$\Z$-分次流形”，使得把 $\Z$-分次换成相应的 $\Z/2\Z$-分次后得到的超流形是 $\Pi T({}^L G/{}^L B)$，即 ${}^L G/{}^L B$ 的切丛，其纤维的奇偶性由偶改为奇。

我们期待范畴 $\hat{\fg}_{\kappa_c}\text{-}\mathrm{mod}_{\mathrm{nilp}}^{I^0}$ 在 $\Op^{\mathrm{nilp}}_{{}^L G}$ 上是平坦的。因此把 猜想 10.4.4 特殊化到特定的 oper $\chi\in\Op^{\mathrm{nilp}}_{{}^L G}$，得到如下推论：范畴等价
\begin{equation}
    D^b\bigl(\hat{\fg}_{\kappa_c}\text{-}\mathrm{mod}_\chi\bigr)^{I^0}
    \simeq D^b\bigl(\mathrm{QCoh}(\widetilde{\operatorname{Sp}}^{\mathrm{DG}}_{\Res_{\mathcal{F}}(\chi)})\bigr).
    \tag{10.4-14}
\end{equation}
这与 猜想 10.1.2 相符：该猜想说若 opers $\chi_1$ 与 $\chi_2$ 的底层局部系同构，则范畴 $\hat{\fg}_{\kappa_c}\text{-}\mathrm{mod}_{\chi_1}$ 与 $\hat{\fg}_{\kappa_c}\text{-}\mathrm{mod}_{\chi_2}$（从而 $D^b(\hat{\fg}_{\kappa_c}\text{-}\mathrm{mod}_{\chi_1})^{I^0}$ 与 $D^b(\hat{\fg}_{\kappa_c}\text{-}\mathrm{mod}_{\chi_2})^{I^0}$）应当等价。对幂零 opers $\chi_1,\chi_2$，这当且仅当它们的单值性彼此共轭。由于它们的单值性由留数取指数得到，这等价于说留数 $\Res_{\mathcal{F}}(\chi_1)$ 与 $\Res_{\mathcal{F}}(\chi_2)$ 关于 ${}^L G$ 的 $\mathcal{F}_{{}^L B,0}$-扭变彼此共轭。但此时 $\chi_1$ 与 $\chi_2$ 对应的 DG Springer 纤维也同构，故由 (10.4-14)，$D^b(\hat{\fg}_{\kappa_c}\text{-}\mathrm{mod}_{\chi_1})^{I^0}$ 与 $D^b(\hat{\fg}_{\kappa_c}\text{-}\mathrm{mod}_{\chi_2})^{I^0}$ 彼此等价。

范畴 $D^b(\mathrm{QCoh}(\widetilde{\operatorname{Sp}}^{\mathrm{DG}}_u))$（$u$ 为任意幂零元）的 Grothendieck 群，与 $\mathrm{QCoh}(\operatorname{Sp}_u)$ 的 Grothendieck 群相同。换言之，Grothendieck 群“不知道”$\widetilde{\operatorname{Sp}}^{\mathrm{DG}}_u$ 的 DG 结构或非约化结构。因此它无非就是代数 $K$-理论 $K(\operatorname{Sp}_u)$。如我们在 10.4.1 节末尾所解释，这个代数 $K$-理论的等变变体实现了仿射 Hecke 代数 $H(G(F),I)$ 上的“标准模”。此外，上面那样的 $I$-不变向量空间 $R^{\pi(I)}$（它们自然是仿射 Hecke 代数上的模）可以实现为 $K(\operatorname{Sp}_u)$ 的子商。这表明等价 (10.4-14) 与 (10.4-13) 和经典结果相容。

% [Tutor Note: p.350 易错处。$K_0(D^b(\mathrm{QCoh}(\widetilde{\operatorname{Sp}}^{\mathrm{DG}}_u)))=K(\operatorname{Sp}_u)$：DG 与非约化结构只影响高阶信息，Grothendieck 群只看 $K_0$。这正是为什么用导出范畴能“容纳”经典 $K$-理论内容；若只在阿贝尔范畴层面比较，会丢失这些信息。]
然而乍看之下，经典图像与范畴图像之间有一些重要差别，下面更详细地讨论。

在 10.4.1 节概述的 $H(G(F),I)$-模构造中，我们须选取 ${}^L G$ 的半单元素 $\gamma$，使 $\gamma u\gamma^{-1}=qu$，其中 $q$ 是 $F$ 的剩余域的元素个数。然后考虑特殊化的 $A$-等变 $K$-理论 $K^A(\operatorname{Sp}_u)_{(\gamma,q)}$，其中 $A$ 是 ${}^L G\times\C^\times$ 中含 $(\gamma,q)$ 的最小子代数群。这给出 $K(\operatorname{Sp}_u)$ 上的 $H(G(F),I)$-模结构。但这个模带有关于群 $C(\gamma,u)$（$\gamma$ 与 $u$ 在 ${}^L G$ 中中心化子的分量群）的剩余对称性，且该对称性与 $H(G(F),I)$ 的作用交换。因此考虑对应于 $C(\gamma,u)$ 的不可约表示 $\rho$ 的 $H(G(F),I)$-模
\[
    K^A(\operatorname{Sp}_u)_{(\gamma,q,\rho)}
    =\operatorname{Hom}_{C(\gamma,u)}\bigl(\rho,K(\operatorname{Sp}_u)\bigr).
\]
最后，这些分量每个都有唯一的不可约商，这就是 $H(G(F),I)$ 的不可约表示，它实现在空间 $R^{\pi(I)}$ 上，其中 $(R,\pi)$ 是在双射 (10.4-6) 下对应于 $(\gamma,u,\rho)$ 的 $G(F)$ 的不可约表示。这一精细结构如何反映在范畴情形中？

我们的范畴 $D^b(\mathrm{QCoh}(\widetilde{\operatorname{Sp}}^{\mathrm{DG}}_u))$（其中 $u=\Res_{\mathcal{F}}(\chi)$）是（非等变）$K$-理论 $K(\operatorname{Sp}_u)$ 的一种特定范畴化。注意在经典局部 Langlands 对应 (10.4-6) 中，三元组 $(\gamma,u,\rho)$ 的元素 $u$ 被解释为相应 Weil--Deligne 群 $W'_F$ 表示的单值性的对数。这与把 $\Res_{\mathcal{F}}(\chi)$ 解释为对应于 $\chi$ 的 $D^\times$ 上 ${}^L G$-局部系（它扮演范畴 $\hat{\fg}_{\kappa_c}\text{-}\mathrm{mod}_\chi$ 的局部 Langlands 参数的角色）单值性的对数一致（相差不重要因子 $2\pi i$）。\transnote{经典参数 $(\gamma,u,\rho)$ 中只有 $u$ 有直接的几何对应：$u$ 是留数 $\operatorname{Res}_{\mathcal F}(\chi)$（相差无关紧要的 $2\pi i$），其单值性幂幺正是“留数幂零”的几何表现。$\gamma$（Frobenius）与 $\rho$（分量群表示）都没有直接类比物：几何情形只固定一个局部系 $\sigma$，其自同构群 $\operatorname{Aut}(\sigma)$ 取代了 $\gamma$ 的位置。这也是 (10.4-14) 右端取非等变拟凝聚层范畴、而无分歧情形取带 ${}^L G$ 作用的等变版本的原因。[来源：本书 p.350–352；[Frenkel and Gaitsgory 2006c]（本书 p.351 引用）。]}

但另外两个参数 $\gamma$ 与 $\rho$ 呢？为什么我们的范畴对应 Springer 纤维的非等变 $K$-理论，而不是像经典情形那样的等变 $K$-理论？

对应于 $W'_F$ 中 Frobenius 的元素 $\gamma$，在几何情形似乎没有类比物。我们在上面的无分歧情形已看到这一点：在经典情形无分歧局部 Langlands 参数是 ${}^L G$ 中的半单共轭类 $\gamma$，而在几何情形我们只有一个无分歧局部 Langlands 参数，即平凡局部系。

为更好地理解这里发生的事，重新考察无分歧情形。回忆球面 Hecke 代数 $H(G(F),K_0)$ 同构于表示环 $\Rep{}^L G$。不可约无分歧表示 $(R,\pi)$ 的一维 $K_0$-不变空间实现 $H(G(F),K_0)$ 的一维表示，即同态 $\Rep{}^L G\to\C$。$(R,\pi)$ 的无分歧 Langlands 参数 $\gamma$（${}^L G$ 中的半单共轭类）是对应于该同态的点 $\Spec(\Rep{}^L G)$。这个同态的范畴类比物是什么？$\Rep{}^L G$ 的范畴化是范畴 $\mathcal{R}ep{}^L G$。$\Rep{}^L G$ 的乘积结构反映在 $\mathcal{R}ep{}^L G$ 的张量范畴结构上。另一方面，代数 $\C$ 的范畴化是向量空间范畴 $\mathcal{V}ect$。因此同态 $\Rep{}^L G\to\C$ 的范畴类比物，是保持两个范畴张量结构的函子 $\mathcal{R}ep{}^L G\to\mathcal{V}ect$。这样的函子称为纤维函子。纤维函子自身构成范畴，它等价于 ${}^L G$-torsor 的范畴。因此任意两个纤维函子同构，但不同构典范地。特别地，每个纤维函子的自同构群同构于 ${}^L G$。（顺便说，Tannakian 形式中正是这样从纤维函子重构 ${}^L G$。）因此我们看到，在范畴世界里我们虽没有半单共轭类 $\gamma$（即 $\Spec(\Rep{}^L G)$ 的点）的类比物，它们的角色在某种意义上由纤维函子的自同构群扮演。

% [Tutor Note: p.351--352。这段是本章最富概念性的段落之一：经典参数 $\gamma$ 的“位置”被“纤维函子的自同构群”取代，而纤维函子的范畴等价于 ${}^L G$-torsor 的范畴。于是“选一个 $\gamma$”在几何上对应“选一个纤维函子”（即选一个 torsor）；经典情形只有一个无分歧参数，几何情形只有一个平凡局部系，但其自同构群 ${}^L G$ 提供了原本由 $\gamma$ 连续变化提供的参数空间。]
这反映在如下事实中：虽然在范畴情形我们只有一个无分歧 Langlands 参数——即 $D^\times$ 上的平凡 ${}^L G$-局部系 $\sigma_0$——但这个局部系有非平凡的自同构群，即 ${}^L G$。因此我们期待群 ${}^L G$ 通过自同构作用在对应于 $\sigma_0$ 的 Langlands 范畴 $\mathcal{C}_{\sigma_0}$ 上，且这个作用应与圈群 $G((t))$ 在 $\mathcal{C}_{\sigma_0}$ 上的作用交换。正是 ${}^L G$ 的这个作用，意在弥补与经典情形相比缺少的无分歧 Langlands 参数。

我们在 10.3 节论证过，范畴 $\hat{\fg}_{\kappa_c}\text{-}\mathrm{mod}_\chi$（其中
\[
    \chi=(\mathcal{F},\nabla,\mathcal{F}_{{}^L B})\in\Op_{{}^L G}(D),
\]
）是 Langlands 范畴 $\mathcal{C}_{\sigma_0}$ 的候选。因此我们期待群 ${}^L G$（更精确地说，其扭变 ${}^L G_{\mathcal{F}}$）作用在范畴 $\hat{\fg}_{\kappa_c}\text{-}\mathrm{mod}_\chi$ 上。在 10.3.7 节，我们展示了如何用范畴 $\hat{\fg}_{\kappa_c}\text{-}\mathrm{mod}_\chi$ 与仿射 Grassmann 流形 $\Gr$ 上 Hecke 本征模的范畴 $\mathcal{D}^{\mathrm{Hecke}}_{\kappa_c}\text{-}\mathrm{mod}_\chi$ 之间的猜想等价（见 猜想 10.3.9）得到这个作用。范畴 $\mathcal{D}^{\mathrm{Hecke}}_{\kappa_c}\text{-}\mathrm{mod}_\chi$ 在 10.3.6 节中定义为 $\Gr$ 上扭变 D-模的范畴 $\mathcal{D}_{\kappa_c}\text{-}\mathrm{mod}$ 关于范畴 $\mathcal{R}ep{}^L G$ 的幺半作用的“去等变化”。\transnote{$\mathcal D^{\mathrm{Hecke}}_{\kappa_c}\text{-}\mathrm{mod}_\chi$ 是整体几何 Langlands 在仿射 Grassmann 流形一侧的 Hecke 本征条件范畴（即 Hecke 本征模，eigenmodule，不要与整体情形的本征层 eigensheaf 混淆）：Hecke 算子（来自 $\mathcal H(G((t)),G[[t]])\simeq\Rep{}^L G$ 的每个对象）在它上面按本征值 $\chi$ 作用。这里的“去等变化”是 Arkhipov–Gaitsgory 的通用操作，与“等变化”互为逆：先保留 $\Rep{}^L G$ 的幺半作用，再把它剥掉，从而得到带 ${}^L G$ 作用的范畴；它绝不是取 $G[[t]]$-等变对象。10.3.5 节的“Hecke 本征模”是本征条件的模侧表述，本段则是它在 $\Gr$ 一侧的几何表述，两者层次不同、不可混用。[来源：本书 p.352；[Arkhipov and Gaitsgory 2003]（本书 p.331 引用）；[Frenkel and Gaitsgory 2006c]（本书 p.352 引用）。]}

现在来到一个对理解温和分歧情形运作方式很有用的关键观察：范畴 $\mathcal{R}ep{}^L G$ 可以解释为簇 $\mathrm{pt}=\Spec\C$ 上 ${}^L G$-等变拟凝聚层的范畴。换言之，$\mathcal{R}ep{}^L G$ 可以解释为栈 $\mathrm{pt}/{}^L G$ 上拟凝聚层的范畴。范畴 $\mathcal{R}ep{}^L G$ 在 $\mathcal{D}_{\kappa_c}\text{-}\mathrm{mod}$ 上存在幺半作用，应当看作“范畴 $\mathcal{D}_{\kappa_c}\text{-}\mathrm{mod}$ 纤维化于栈 $\mathrm{pt}/{}^L G$ 之上”的陈述。于是 猜想 10.3.9 可以解释为说
\[
    \hat{\fg}_{\kappa_c}\text{-}\mathrm{mod}_\chi
    \simeq\mathcal{D}_{\kappa_c}\text{-}\mathrm{mod}
    \times_{\mathrm{pt}/{}^L G}\mathrm{pt}.
\]
换言之，若 $\mathcal{C}$ 是纤维化于 $D^\times$ 上所有 ${}^L G$-局部系的栈 $\Loc_{{}^L G}(D^\times)$ 之上的猜想 Langlands 范畴，则
\[
    \mathcal{D}_{\kappa_c}\text{-}\mathrm{mod}
    \simeq\mathcal{C}\times_{\Loc_{{}^L G}(D^\times)}\mathrm{pt}/{}^L G,
\]
而
\[
    \hat{\fg}_{\kappa_c}\text{-}\mathrm{mod}_\chi
    \simeq\mathcal{C}\times_{\Loc_{{}^L G}(D^\times)}\mathrm{pt},
\]
其中态射 $\mathrm{pt}\to\Loc_{{}^L G}(D^\times)$ 对应于 oper $\chi$。

% [Tutor Note: p.352 易错处。这里出现了“点”的两种解释：平凡局部系 $\sigma_0$ 既可以视为一个点（由具体的 ${}^L G$-丛加联络给出，如正则 oper $\chi$），也可以视为栈 $\mathrm{pt}/{}^L G$（因为平凡局部系有自同构群 ${}^L G$）。前者给出带 ${}^L G$ 作用的范畴，后者给出带 $\mathcal{R}ep{}^L G$ 幺半作用的范畴；两者通过等变化/去等变化互相转换。这是理解“为什么几何情形要多带一个对称群作用”的关键。]
因此在范畴情形，关于平凡局部系 $\sigma_0$ 有两种不同看法：既可以看作一个点（由 $D$ 上某个具体的 ${}^L G$-丛加联络给出，例如正则 oper $\chi$），也可以看作栈 $\mathrm{pt}/{}^L G$。第一种情形下 Langlands 范畴的基变换给出带 ${}^L G$ 作用的范畴，例如 $\hat{\fg}_{\kappa_c}\text{-}\mathrm{mod}_\chi$ 或 $\mathcal{D}^{\mathrm{Hecke}}_{\kappa_c}\text{-}\mathrm{mod}$。第二种情形下的基变换给出带 $\mathcal{R}ep{}^L G$ 幺半作用的范畴，例如 $\mathcal{D}_{\kappa_c}\text{-}\mathrm{mod}$。通过分别关于 ${}^L G$ 与 $\mathcal{R}ep{}^L G$ 施行等变化与去等变化，我们可以在两个版本之间来回转换。

现在回到温和分歧情形。三元组 $(\gamma,u,\rho)$ 中出现的半单元素 $\gamma$，扮演与无分歧 Langlands 参数 $\gamma$ 相同的角色。但现在它须满足恒等式 $\gamma u\gamma^{-1}=qu$。回忆 $H(G(F),I)$ 的中心 $Z$ 同构于 $\Rep{}^L G$，因此 $\Spec Z$ 是 ${}^L G$ 中所有半单元素的集合。对固定的幂零元 $u$，方程 $\gamma u\gamma^{-1}=qu$ 在 $\Spec Z$ 中截出一个轨迹 $C_u$，对应那些可能出现在对应于 $u$ 的不可约 $H(G(F),I)$-模上的中心特征标。在范畴情形（此时令 $q=1$），$C_u$ 的类比物是 $u$ 在 ${}^L G$ 中的中心化子 $Z(u)$，它恰是 $D^\times$ 上单值性为 $\exp(2\pi iu)$ 的温和分歧局部系 $\sigma$ 的自同构群 $\operatorname{Aut}(\sigma)$。按一般原理，我们期待群 $\operatorname{Aut}(\sigma)$ 作用在 Langlands 范畴 $\mathcal{C}_\sigma$ 上，正如我们期待平凡局部系 $\sigma_0$ 的自同构群 ${}^L G$ 作用在范畴 $\mathcal{C}_{\sigma_0}$ 上。正是这个作用取代了几何情形中的参数 $\gamma$。\transnote{固定 $u$ 后，方程 $\gamma u\gamma^{-1}=qu$ 在 $\Spec Z=({}^L G)^{\mathrm{ss}}$ 中截出的 $C_u$ 记录“哪些 $\gamma$ 与 $u$ 相容”。几何情形令 $q=1$，$C_u$ 退化为 $u$ 的中心化子 $Z(u)$；而 $Z(u)$ 正是单值性为 $\exp(2\pi iu)$ 的局部系 $\sigma$ 的自同构群。于是“选一个 Frobenius 参数”在几何上变成“选一个局部系并带上其自同构群的作用”，这与无分歧情形用 $\sigma_0$ 的自同构群 ${}^L G$ 补足参数是同一机理。[来源：本书 p.353；[Frenkel and Gaitsgory 2006c]（本书 p.352 引用）。]}

在经典情形我们还有一个参数 $\rho$。回忆 $\rho$ 是中心化子 $Z(\gamma,u)$（$\gamma$ 与 $u$ 的共同中心化子）的分量群 $C(\gamma,u)$ 的表示。但群 $Z(\gamma,u)$ 是群 $Z(u)$ 的子群，而后者在几何情形成为群 $\operatorname{Aut}(\sigma)$。因此可以说参数 $\rho$ 也被吸收进 $\operatorname{Aut}(\sigma)$ 在范畴 $\mathcal{C}_\sigma$（或在其众多化身 $\hat{\fg}_{\kappa_c}\text{-}\mathrm{mod}_\chi$，$\chi\in\Op^{\mathrm{nilp}}_{{}^L G}$ 之一）上的作用。

% [Tutor Note: p.353。这里解释了“为什么 (10.4-14) 右端是非等变层范畴”：因为自同构群 $\operatorname{Aut}(\sigma)$ 的作用必须先“去等变化”（de-equivariantize），才能与经典情形中有 $\rho$ 分量的 $K$-理论对应。与无分歧情形（那里用带 ${}^L G$ 作用的范畴）相比，这里最终取的是去等变化后的版本，所以 Grothendieck 群对应非等变 $K$-理论。]
若我们有 $\operatorname{Aut}(\sigma)$ 在范畴 $\mathcal{C}_\sigma$（或其化身 $\hat{\fg}_{\kappa_c}\text{-}\mathrm{mod}_\chi$，$\chi\in\Op^{\mathrm{nilp}}_{{}^L G}$）上的作用，这意味着这些范畴必须“去等变化”，就像无分歧情形中范畴 $\hat{\fg}_{\kappa_c}\text{-}\mathrm{mod}_\chi$（$\chi\in\Op_{{}^L G}(D)$）那样。这就是为什么在等价 (10.4-14)（以及 猜想 10.4.4）中出现的是非等变的拟凝聚层范畴（其 Grothendieck 群对应 Springer 纤维的非等变 $K$-理论）。

然而这些范畴也有等变版本。考虑 10.4.2 节引入的 $\Loc_{{}^L G}(D^\times)$ 中温和分歧局部系的子栈 $\Loc^{\mathrm{RS,tame}}_{{}^L G}$。由于温和分歧局部系完全由其（幂幺）单值性的对数决定，这个子栈同构于 $\mathcal{N}/{}^L G$。这个子栈扮演对应于平凡局部系的子栈 $\mathrm{pt}/{}^L G$ 的角色。令
\[
    \mathcal{C}_{\mathrm{tame}}
    =\mathcal{C}\times_{\Loc_{{}^L G}(D^\times)}\mathcal{N}/{}^L G.
\]
于是按 (10.1-3) 的笛卡尔方块所表达的一般猜想，我们期待
\begin{equation}
    \hat{\fg}_{\kappa_c}\text{-}\mathrm{mod}_{\mathrm{nilp}}
    \simeq\mathcal{C}_{\mathrm{tame}}
    \times_{\mathcal{N}/{}^L G}\Op^{\mathrm{nilp}}_{{}^L G}.
    \tag{10.4-15}
\end{equation}
设 $D^b(\mathcal{C}_{\mathrm{tame}})^{I^0}$ 是对应于 $\mathcal{C}_{\mathrm{tame}}$ 的 $I^0$-等变导出范畴。把 (10.4-15) 与 猜想 10.4.4 结合，并注意
\[
    \widetilde{\mathrm{MOp}}{}^0_{{}^L G}
    \simeq\Op^{\mathrm{nilp}}_{{}^L G}
    \times_{\mathcal{N}/{}^L G}
    {}^{L}\widetilde{\widetilde{\mathcal{N}}}/{}^L G,
\]
我们得到如下猜想（见 [Frenkel and Gaitsgory 2006c]）：
\begin{equation}
    D^b(\mathcal{C}_{\mathrm{tame}})^{I^0}
    \simeq D^b\bigl(\mathrm{QCoh}({}^{L}\widetilde{\widetilde{\mathcal{N}}}/{}^L G)\bigr).
    \tag{10.4-16}
\end{equation}
右端的范畴可以解释为“加厚”Springer 分解 ${}^L\widetilde{\widetilde{\mathcal{N}}}$ 上 ${}^L G$-等变拟凝聚层的导出范畴。猜想性等价 (10.4-14) 与 (10.4-16) 合起来，应当看作仿射 Hecke 代数的模在 Springer 纤维的 $K$-理论中的实现这一经典结果的范畴版本。

等价 (10.4-14) 的一个推论如下：范畴 $\hat{\fg}_{\kappa_c}\text{-}\mathrm{mod}_\chi^{I^0}$ 的不可约对象在 $\hat{\fg}_{\kappa_c}\text{-}\mathrm{mod}_\chi^{I^0}$ 的 Grothendieck 群中的类，给出代数 $K$-理论 $K(\operatorname{Sp}_u)$ 的一组基，其中 $u=\Res_{\mathcal{F}}(\chi)$。据推测，这组基与 G. Lusztig 在 [1998] 中构造的（此 $K$-理论等变版本的）基密切相关（后者是从正特征下无限制 $\fg$-模的角度出发的）。

\subsection{猜想的证据}

\begin{keypoints}
    \item \textbf{核心思想：} 主猜想有四组证据：Wakimoto 模对应摩天楼层；某些商范畴的等价（定理 10.4.5）；正则 opers 的截断等价；与 Bezrukavnikov 理论的联系。
    \item \textbf{数学工具：} Wakimoto 模 $W_{\overline{\nabla}}$ 与 Miura 变换 $\mu$；Springer 纤维的层化；$\mu^{-1}(\Op^{\mathrm{nilp}})$ 的分层结构；范畴 Hecke 代数在扭变 D-模范畴上的作用；中心层。
    \item \textbf{关键结论：} 定理 10.4.5 证明（在阿贝尔商范畴层面）猜测的一部分；正则 opers 的截断给出 (10.4-17) 与 (10.4-13) 的又一证据；Bezrukavnikov 的等价 (10.4-18) 给出与仿射旗簇一侧的对应。
    \item \textbf{记号提醒：} $\mu$ 是 Miura 变换；$\mathrm{Conn}(\Omega^{-\rho})$ 是联络空间；$\mathfrak{sl}_2^{(i)}$ 是 $\fg$ 中由 $e_i,h_i,f_i$ 生成的子代数；$\mathrm{Fl}=G((t))/I$。
\end{keypoints}

% [Tutor Note: p.354。本小节是“证据汇编”，四组证据分别从不同角度支持 猜想 10.4.4：(1) 具体对象层面（Wakimoto 模 ↔ 摩天楼层）；(2) 范畴层面（商范畴等价，定理 10.4.5）；(3) 截断到正则 opers；(4) 与 Bezrukavnikov 的 $K$-理论/导出等价理论对接。读的时候把握“层—模”对应这条主线即可。]
现在描述 猜想 10.4.4 的一些证据。它由以下四组结果组成：
\begin{itemize}
    \item 把 Wakimoto 模解释为对应于 $\widetilde{\mathrm{MOp}}{}^0_{{}^L G}$ 上摩天楼层的 $\hat{\fg}_{\kappa_c}$-模。
    \item 证明 $D^b(\hat{\fg}_{\kappa_c}\text{-}\mathrm{mod}_{\mathrm{nilp}})^{I^0}$ 与 $D^b(\mathrm{QCoh}(\widetilde{\mathrm{MOp}}{}^0_{{}^L G}))$ 的某些商范畴彼此等价。
    \item 证明等价 (10.4-13) 到正则 opers 的限制。
    \item 与 R. Bezrukavnikov 理论的联系。
\end{itemize}
我们从 Wakimoto 模的讨论开始。

假设我们已证明范畴等价 (10.4-13)。则 $\widetilde{\mathrm{MOp}}{}^0_{{}^L G}$ 上每个拟凝聚层应对应导出范畴 $D^b(\hat{\fg}_{\kappa_c}\text{-}\mathrm{mod}_{\mathrm{nilp}})^{I^0}$ 的一个对象。$\widetilde{\mathrm{MOp}}{}^0_{{}^L G}$ 上最简单的拟凝聚层是支撑在 $\widetilde{\mathrm{MOp}}{}^0_{{}^L G}$ 的 $\C$-点上的摩天楼层。由定义，$\widetilde{\mathrm{MOp}}{}^0_{{}^L G}$ 的 $\C$-点（等同于约化概形 $\mathrm{MOp}^0_{{}^L G}$ 的 $\C$-点）是一对 $(\chi,\mathfrak{b}')$，其中 $\chi=(\mathcal{F},\nabla,\mathcal{F}_{{}^L B})$ 是 $\Op^{\mathrm{nilp}}_{{}^L G}$ 中的幂零 ${}^L G$-oper，$\mathfrak{b}'$ 是对应于 $\Res_{\mathcal{F}}(\chi)$ 的 Springer 纤维的一个点，即 ${}^L\fg_{\mathcal{F}_0}$ 中含 $\Res_{\mathcal{F}}(\chi)$ 的 Borel 子代数构成的簇的一个点。因此，若 猜想 10.4.4 成立，我们应当有一族由这些数据参数化的范畴 $D^b(\hat{\fg}_{\kappa_c}\text{-}\mathrm{mod}_{\mathrm{nilp}})^{I^0}$ 的对象。这些对象是什么？

研读过本书上一章的读者无疑已经猜到答案：它们是 Wakimoto 模！确实，回忆临界水平的 Wakimoto 模由空间 $\mathrm{Conn}(\Omega^{-\rho})_{D^\times}$ 参数化。对应于 $\overline{\nabla}\in\mathrm{Conn}(\Omega^{-\rho})_{D^\times}$ 的 Wakimoto 模记作 $W_{\overline{\nabla}}$（见 9.4.3 节）。根据定理 8.3.3，中心 $Z(\hat{\fg})$ 通过中心特征标 $\mu(\overline{\nabla})$ 作用在 $W_{\overline{\nabla}}$ 上，其中 $\mu$ 是 Miura 变换。此外，如 9.4.3 节所解释，若 $\chi\in\Op^{\mathrm{nilp}}_{{}^L G}$，则对任意 $\overline{\nabla}\in\mu^{-1}(\chi)$，$W_{\overline{\nabla}}$ 是范畴 $\hat{\fg}_{\kappa_c}\text{-}\mathrm{mod}_\chi^I$ 的对象。

根据定理 9.3.7，$\mu^{-1}(\chi)$ 的点与对应于幂零元 $\Res_{\mathcal{F}}(\chi)$ 的 Springer 纤维 $\operatorname{Sp}_{\Res_{\mathcal{F}}(\chi)}$ 的点一一对应。因此，对 $\operatorname{Sp}_{\Res_{\mathcal{F}}(\chi)}$ 的每个点，我们配上一个 Wakimoto 模，它是范畴 $\hat{\fg}_{\kappa_c}\text{-}\mathrm{mod}_\chi^{I^0}$（从而是相应导出范畴）的对象。换言之，Wakimoto 模是范畴 $\hat{\fg}_{\kappa_c}\text{-}\mathrm{mod}_{\mathrm{nilp}}^{I^0}$ 中由 $\widetilde{\mathrm{MOp}}{}^0_{{}^L G}$ 的 $\C$-点参数化的一族对象。自然地假设它们在等价 (10.4-13) 下对应于 $\widetilde{\mathrm{MOp}}{}^0_{{}^L G}$ 上的摩天楼层。这实际上是我们提出这个猜想的动机之一。\transnote{这是 §10.4.5 四组证据中的第一组，也是唯一直接给出“点—对象”对应的一组：Wakimoto 模是最具体的可计算对象（自由场实现），摩天楼层是最具体的拟凝聚层，而 Miura 变换把它们粘起来——定理 8.3.3 使中心按 $\mu(\overline\nabla)$ 作用，定理 9.3.7 使 $\mu^{-1}(\chi)$ 与 Springer 纤维同构。正因为每个 $(\chi,\mathfrak b')$ 点都能落实为一个 Wakimoto 模，作者才提出 (10.4-13) 这一整体陈述。[来源：本书 p.354–355；[Frenkel and Gaitsgory 2006c]（本书 p.355 引用）。]}

% [Tutor Note: p.355 背景。Wakimoto 模 ↔ 摩天楼层是本书“自由场实现”与几何 Langlands 的接点：Wakimoto 模是最具体的可计算对象，摩天楼层是最具体的拟凝聚层，二者通过 Miura 变换联系起来。注意这里说“对每个点配上一个模”，即 Morita/层-模范畴层面的对应，而不是单个模的函子性。]
顺便说，这让我们得以一窥 oper $\chi$ 底层 ${}^L G$-局部系的自同构群如何作用在范畴 $\hat{\fg}_{\kappa_c}\text{-}\mathrm{mod}_\chi$ 上。这个群是 $Z(\Res_{\mathcal{F}}(\chi))$，即留数 $\Res_{\mathcal{F}}(\chi)$ 的中心化子，它作用在 Springer 纤维 $\operatorname{Sp}_{\Res_{\mathcal{F}}(\chi)}$ 上。因此 $g\in Z(\Res_{\mathcal{F}}(\chi))$ 把支撑在点 $p\in\operatorname{Sp}_{\Res_{\mathcal{F}}(\chi)}$ 的摩天楼层送到支撑在 $g\cdot p$ 的摩天楼层。于是我们期待 $g$ 把对应于 $p$ 的 Wakimoto 模送到对应于 $g\cdot p$ 的 Wakimoto 模。

若 Wakimoto 模确实对应摩天楼层，则等价 (10.4-13) 可以看作范畴 $D^b(\hat{\fg}_{\kappa_c}\text{-}\mathrm{mod}_{\mathrm{nilp}})^{I^0}$ 的一种“谱分解”，其基本对象是 Wakimoto 模 $W_{\overline{\nabla}}$，其中 $\overline{\nabla}$ 取遍 $\mathrm{Conn}(\Omega^{-\rho})_{D^\times}$ 中的轨迹，该轨迹逐点同构于 $\widetilde{\mathrm{MOp}}{}^0_{{}^L G}$。

然而要注意：虽然 $\Op^{\mathrm{nilp}}_{{}^L G}$ 在 $\mathrm{Conn}(\Omega^{-\rho})_{D^\times}$ 中在 Miura 变换 $\mu$ 下的（集合论）原像等于 $\mathrm{MOp}^0_{{}^L G}$，但作为代数簇它们并不同构。相反，根据 9.3.3 节，$\mu^{-1}(\Op^{\mathrm{nilp}}_{{}^L G})$ 是子簇的不交并
\[
    \mu^{-1}\bigl(\Op^{\mathrm{nilp}}_{{}^L G}\bigr)
    =\bigsqcup_{w\in W}\mathrm{Conn}(\Omega^{-\rho})^{w(\rho)-\rho}_D,
\]
其中 $\mathrm{Conn}(\Omega^{-\rho})^{w(\rho)-\rho}_D$ 是具正则奇点且留数为 $w(\rho)-\rho$ 的联络构成的簇。另一方面，$\mathrm{MOp}^0_{{}^L G}$ 有一个分层（由与旗簇 ${}^L G/{}^L B$ 上 Schubert 胞腔的交得到），其层片 $\mathrm{MOp}^{0,w}_{{}^L G}$ 同构于 $\mathrm{Conn}(\Omega^{-\rho})^{w(\rho)-\rho}_D$，$w\in W$（见 9.3.2 节与定理 9.3.6）。但这些层片在 $\mathrm{MOp}^0_{{}^L G}$ 中以非平凡的方式“粘合”在一起。

对每个层片（它是仿射代数簇，事实上是无穷维仿射空间），可以用公式 9.4.3 定义的 Wakimoto 模 $W_{w(\rho)-\rho}$ 配上一个函子
\[
    F_w:\mathrm{QCoh}\bigl(\mathrm{Conn}(\Omega^{-\rho})^{w(\rho)-\rho}_D\bigr)
    \to\hat{\fg}_{\kappa_c}\text{-}\mathrm{mod}_{\mathrm{nilp}}^{I^0}.
\]
函子 $F_w$ 把 $\mathrm{QCoh}(\mathrm{Conn}(\Omega^{-\rho})^{w(\rho)-\rho}_D)$ 的对象（等同于代数 $\operatorname{Fun}\mathrm{Conn}(\Omega^{-\rho})^{w(\rho)-\rho}_D$ 上的模 $\mathcal{F}$）送到 $\hat{\fg}_{\kappa_c}$-模
\[
    W_{w(\rho)-\rho}
    \underset{\operatorname{Fun}\mathrm{Conn}(\Omega^{-\rho})^{w(\rho)-\rho}_D}{\otimes}
    \mathcal{F}.
\]
非正式地说，范畴 $D^b(\mathrm{QCoh}(\widetilde{\mathrm{MOp}}{}^0_{{}^L G}))$ 是由各层片 $\mathrm{Conn}(\Omega^{-\rho})^{w(\rho)-\rho}_D$ 上拟凝聚层范畴“粘合”而成的。于是 猜想 10.4.4 可以解释为：范畴 $D^b(\hat{\fg}_{\kappa_c}\text{-}\mathrm{mod}_{\mathrm{nilp}})^{I^0}$ 是由这些函子的像所得到的子范畴“粘合”而成。详见 [Frenkel and Gaitsgory 2006b]。

% [Tutor Note: p.356。这里出现两处“粘合”：几何一侧 $\mathrm{MOp}^0$ 的层片沿 Schubert 胞腔粘合；代数一侧 $\hat{\fg}_{\kappa_c}$-模范畴由各层片上的函子 $F_w$ 的像粘合。猜想 10.4.4 的实质正是这两种粘合方式一致。注意 $\mathrm{MOp}^0$ 的层片是无穷维仿射空间，这与 $\mathrm{Conn}(\Omega^{-\rho})$ 的仿射结构相符。]
此时把猜想性等价 (10.4-13) 与正则 opers 情形的等价 (10.3-9) 作比较是有启发的。把后者改写为更有提示性的形式：
\begin{equation}
    \hat{\fg}_{\kappa_c}\text{-}\mathrm{mod}_{\mathrm{reg}}^{G[[t]]}
    \simeq\mathrm{QCoh}\bigl(\Op_{{}^L G}(D)\bigr).
    \tag{10.4-17}
\end{equation}
这两个陈述之间有两个本质区别。首先，(10.4-17) 是阿贝尔范畴之间的等价，而在 (10.4-13) 中须使用导出范畴。第二个区别是：(10.4-17) 中是光滑仿射代数簇 $\Op_{{}^L G}(D)$（实际上是无穷维仿射空间）上拟凝聚层的范畴，而 (10.4-13) 中是高度奇异且非约化的拟射影簇 $\widetilde{\mathrm{MOp}}{}^0_{{}^L G}$ 上拟凝聚层的范畴。在某种意义上，证明 (10.4-13) 的大部分困难是由 $\widetilde{\mathrm{MOp}}{}^0_{{}^L G}$ 的非仿射性造成的。

然而有一个截断版本的等价 (10.4-13)，它涉及 $\widetilde{\mathrm{MOp}}{}^0_{{}^L G}$ 的仿射子概形上拟凝聚层的范畴，以及范畴 $D^b(\hat{\fg}_{\kappa_c}\text{-}\mathrm{mod}_{\mathrm{nilp}})^{I^0}$ 的某个商。下面给出精确陈述。

对每个 $i=1,\dots,\ell$，设 $\mathfrak{sl}_2^{(i)}$ 是由 $e_i,h_i,f_i$ 生成的 $\fg=\fg\otimes1\subset\hat{\fg}_{\kappa_c}$ 的 $\fsl_2$ 子代数。称 $\hat{\fg}_{\kappa_c}$-模 $M$ 为部分可积的，若存在 $i=1,\dots,\ell$，使得 $\mathfrak{sl}_2^{(i)}$ 在 $M$ 上的作用可指数化为相应群 $SL_2$ 的作用。容易看出部分可积对象在 $\hat{\fg}_{\kappa_c}\text{-}\mathrm{mod}_{\mathrm{nilp}}^{I^0}$ 中构成 Serre 子范畴。设 $f\hat{\fg}_{\kappa_c}\text{-}\mathrm{mod}_{\mathrm{nilp}}^{I^0}$ 是 $\hat{\fg}_{\kappa_c}\text{-}\mathrm{mod}_{\mathrm{nilp}}^{I^0}$ 模去部分可积对象子范畴所得的商范畴。用 $fD^b(\hat{\fg}_{\kappa_c}\text{-}\mathrm{mod}_{\mathrm{nilp}})^{I^0}$ 记 $D^b(\hat{\fg}_{\kappa_c}\text{-}\mathrm{mod}_{\mathrm{nilp}})^{I^0}$ 模去“上同调为部分可积”的对象子范畴所得的三角商范畴。

另一方面，考虑 $\widetilde{\mathrm{MOp}}{}^0_{{}^L G}$ 的子概形 $\widetilde{\mathrm{MOp}}{}^{0,w_0}_{{}^L G}$，它是 $\widetilde{\mathrm{MOp}}{}^0_{{}^L G}$ 在映射 $\widetilde{\mathrm{MOp}}{}^0_{{}^L G}\to{}^L\mathfrak{n}$ 下限制到 ${}^L\mathfrak{n}$ 的正则幂零元轨迹所得子概形的闭包。如 10.4.3 节所见，正则幂零元上的非约化 Springer 纤维同构于 $\Spec h_0$，其中 $h_0=\operatorname{Fun}{}^L\fh/(\operatorname{Fun}{}^L\fh)^W_+$ 是维数为 $|W|$ 的 Artin 代数。因此有
\[
    \widetilde{\mathrm{MOp}}{}^{0,w_0}_{{}^L G}
    \simeq\Op^{\mathrm{nilp}}_{{}^L G}\times\Spec h_0.
\]
于是它确实是 $\widetilde{\mathrm{MOp}}{}^0_{{}^L G}$ 的仿射子概形。相应的约化概形事实上是 9.3.2 节引入的层片 $\mathrm{MOp}^{0,w_0}_{{}^L G}\subset\mathrm{MOp}^0_{{}^L G}$。$\widetilde{\mathrm{MOp}}{}^{0,w_0}_{{}^L G}$ 上拟凝聚层的范畴可以实现为 $\widetilde{\mathrm{MOp}}{}^0_{{}^L G}$ 上拟凝聚层范畴的 Serre 商。

下述结果由 [Frenkel and Gaitsgory 2006c] 证明，它表明这两个商范畴彼此等价（且已在阿贝尔范畴层面成立）。

\begin{theorem}\label{thm:u16-quot}
    范畴 $f\hat{\fg}_{\kappa_c}\text{-}\mathrm{mod}_{\mathrm{nilp}}^{I^0}$ 与 $\mathrm{QCoh}(\widetilde{\mathrm{MOp}}{}^{0,w_0}_{{}^L G})$ 彼此等价。
\end{theorem}

% [Tutor Note: p.357--358。定理 10.4.5 的意义：把“最坏”的非仿射性（$\widetilde{\mathrm{MOp}}{}^0$ 的奇点）截断到 $w_0$-层片附近，此时主猜想已可严格证明（阿贝尔范畴层面）。这是主猜想的第二组证据，也提示了完整的 (10.4-13) 应如何借助分层来证。]
这是本节开头所列 猜想 10.4.4 的第二组证据。\transnote{定理 10.4.5 的意义在“截断”：把高度奇异、非仿射的族 $\widetilde{\mathrm{MOp}}{}^0$ 限制到仿射子概形 $\widetilde{\mathrm{MOp}}{}^{0,w_0}\simeq\Op^{\mathrm{nilp}}\times\Spec h_0$（即 $w_0$-层片的闭包），此时两边都取阿贝尔范畴，且等价可严格证明。这提示完整的 (10.4-13) 应沿 $\mathrm{MOp}^0$ 的分层用粘合或导出方法证明，而不是逐点验证。[来源：本书 p.357–358；[Frenkel and Gaitsgory 2006c]（本书 p.357 引用）。]}

现在讨论第三组证据。 (10.4-13) 两边的范畴都纤维化于 $\Op^{\mathrm{nilp}}_{{}^L G}$ 之上；换言之，代数 $\operatorname{Fun}\Op^{\mathrm{nilp}}_{{}^L G}$ 作用在这些范畴的每个对象上。考虑这些范畴到正则 opers 子概形 $\Op_{{}^L G}(D)\subset\Op^{\mathrm{nilp}}_{{}^L G}$ 的限制。相应的阿贝尔范畴定义为使 $\operatorname{Fun}\Op^{\mathrm{nilp}}_{{}^L G}$ 的作用经过 $\operatorname{Fun}\Op_{{}^L G}(D)$ 分解的满子范畴。相应导出范畴的对象是其上同调具有该性质的复形。

(10.4-13) 右端范畴到子概形 $\Op_{{}^L G}(D)$ 的限制，是 DG 概形
\[
    \widetilde{\mathrm{MOp}}{}^{\mathrm{reg}}_{{}^L G}
    =\widetilde{\mathrm{MOp}}{}^0_{{}^L G}
    \overset{L}{\times}_{\Op^{\mathrm{nilp}}_{{}^L G}}\Op_{{}^L G}(D)
\]
上拟凝聚层的导出范畴。相应的张量积函子不是左正合的，于是我们处于 10.4.3 节所描述的情形。因此须取 DG 纤维积。由于任何正则 oper 的留数都等于 $0$，$\widetilde{\mathrm{MOp}}{}^{\mathrm{reg}}_{{}^L G}$ 在每个 oper $\chi\in\Op_{{}^L G}(D)$ 上的纤维同构于 DG Springer 纤维 $\operatorname{Sp}^{\mathrm{DG}}_0$，即 DG 概形 $\Pi T({}^L G/{}^L B)$，如 10.4.3 节所讨论。

另一方面，(10.4-13) 左端范畴到 $\Op_{{}^L G}(D)$ 的限制是范畴 $D^b(\hat{\fg}_{\kappa_c}\text{-}\mathrm{mod}_{\mathrm{reg}})^{I^0}$，即 $\hat{\fg}_{\kappa_c}\text{-}\mathrm{mod}_{\mathrm{reg}}$ 的导出范畴的 $I^0$-等变部分。根据 [Frenkel and Gaitsgory 2005] 中证明的一个定理（它是 猜想 10.3.9 的 $I^0$-等变版本），有范畴等价
\[
    \hat{\fg}_{\kappa_c}\text{-}\mathrm{mod}_{\mathrm{reg}}^{I^0}
    \simeq\mathcal{D}^{\mathrm{Hecke}}_{\kappa_c}\text{-}\mathrm{mod}_{\mathrm{reg}}^{I^0},
\]
因此相应的导出范畴也等价。另一方面，由 [Arkhipov et al. 2004] 的结果，$D^b(\mathcal{D}^{\mathrm{Hecke}}_{\kappa_c}\text{-}\mathrm{mod}_{\mathrm{reg}})^{I^0}$ 的导出范畴等价于 $D^b(\mathrm{QCoh}(\widetilde{\mathrm{MOp}}{}^{\mathrm{reg}}_{{}^L G}))$。因此 (10.4-13) 中出现的范畴到 $\Op_{{}^L G}(D)$ 的限制彼此等价。

% [Tutor Note: p.358。第三组证据的实质是“截断到正则 opers 后主猜想已成立”。关键是把两边都限制到 $\Op_{{}^L G}(D)\subset\Op^{\mathrm{nilp}}_{{}^L G}$：右边得到 DG 纤维积（因非平坦），左边得到 $I^0$-等变的 reg 范畴，而后者由 [Frenkel and Gaitsgory 2005] 与 [Arkhipov et al. 2004] 已知等价于右边的导出范畴。注意这里引用的是 $I^0$-等变版本的 猜想 10.3.9，而非原版。]
最后讨论第四组证据：与 Bezrukavnikov 理论的联系。

为引出它，回忆 10.3.5 节讨论了范畴球面 Hecke 代数 $\mathcal{H}(G((t)),G[[t]])$ 在范畴 $\hat{\fg}_{\kappa_c}\text{-}\mathrm{mod}_\chi$ 上的作用，其中 $\chi$ 是正则 oper。仿射 Hecke 代数 $H(G(F),I)$ 也有范畴类比物。考虑仿射旗簇 $\mathrm{Fl}=G((t))/I$。范畴仿射 Hecke 代数是范畴 $\mathcal{H}(G((t)),I)$，即 $\mathrm{Fl}=G((t))/I$ 上 D-模导出范畴中由上同调为 $I$-等变的复形组成的满子范畴。这个范畴自然作用在导出范畴 $D^b(\hat{\fg}_{\kappa_c}\text{-}\mathrm{mod}_\chi)^I$ 上。这个作用在等价 (10.4-13) 的另一侧对应于什么？

答案由 R. Bezrukavnikov [2006] 的一个定理给出，它可以看作同构 (10.4-8) 的范畴化：
\begin{equation}
    D^b\bigl(\mathcal{D}^{\mathrm{Fl}}_{\kappa_c}\text{-}\mathrm{mod}\bigr)^{I^0}
    \simeq D^b\bigl(\mathrm{QCoh}(\widetilde{\mathrm{St}})\bigr),
    \tag{10.4-18}
\end{equation}
其中 $\mathcal{D}^{\mathrm{Fl}}_{\kappa_c}\text{-}\mathrm{mod}$ 是 $\mathrm{Fl}$ 上扭变 D-模的范畴，$\widetilde{\mathrm{St}}$ 是“加厚”Steinberg 簇
\[
    \widetilde{\mathrm{St}}
    =\widetilde{\mathcal{N}}\times_{\mathcal{N}}\widetilde{\widetilde{\mathcal{N}}}
    =\widetilde{\mathcal{N}}\times_{{}^L\fg}{}^L\widetilde{\fg}.
\]
在道德上，我们期待 (10.4-18) 中的两个范畴以相容的方式作用在 (10.4-13) 中的两个范畴上。然而严格说来，(10.4-18) 的左端是这样作用的：
\[
    D^b\bigl(\hat{\fg}_{\kappa_c}\text{-}\mathrm{mod}_{\mathrm{nilp}}\bigr)^I
    \to D^b\bigl(\hat{\fg}_{\kappa_c}\text{-}\mathrm{mod}_{\mathrm{nilp}}\bigr)^{I^0},
\]
而 (10.4-18) 的右端是这样作用的：
\[
    D^b\bigl(\mathrm{QCoh}(\mathrm{MOp}^0_{{}^L G})\bigr)
    \to D^b\bigl(\mathrm{QCoh}(\widetilde{\mathrm{MOp}}{}^0_{{}^L G})\bigr).
\]
所以需要一个更精确的陈述，可在 [Bezrukavnikov 2006] 第 4.2 节中找到。或者，可以像 [Bezrukavnikov 2006] 那样，考虑 ${}^L G$ 的仿射辫群的相应作用。

这种相容性的一个特殊情形，涉及范畴 $D^b(\mathcal{D}^{\mathrm{Fl}}_{\kappa_c}\text{-}\mathrm{mod})^I$ 的某些特殊对象，即 [Gaitsgory 2001] 引入的中心层。它们对应仿射 Hecke 代数 $H(G(F),I)$ 的中心元。这些中心元在不可约 $H(G(F),I)$-模上、以及在上面讨论的标准模 $K^A(\operatorname{Sp}_u)_{(\gamma,q,\rho)}$ 上作为标量作用。我们已经论证，范畴 $\hat{\fg}_{\kappa_c}\text{-}\mathrm{mod}_\chi^{I^0}$（$\chi\in\Op^{\mathrm{nilp}}_{{}^L G}$）是这些表示的范畴版本。因此自然地期待它的对象关于来自 $D^b(\mathcal{D}^{\mathrm{Fl}}_{\kappa_c}\text{-}\mathrm{mod})^I$ 的中心层的作用是“本征模”（在 10.3.5 节的意义下）。这确实已由 [Frenkel and Gaitsgory 2006a] 证明。\transnote{(10.4-18) 是 Kazhdan–Lusztig 同构 (10.4-8) 的范畴化：左边是 $\mathrm{Fl}$ 上扭变 D-模的导出范畴，右边是加厚 Steinberg 簇 $\widetilde{\mathrm{St}}$ 上的拟凝聚层。中心层（[Gaitsgory 2001]）在左边对应 Hecke 代数的中心元，其作用在不可约模与标准模上为标量；因此自然期待对应的几何对象是“本征模”（eigenmodule），这正是 [Frenkel and Gaitsgory 2006a] 所证明的相容性。[来源：本书 p.359；[Bezrukavnikov 2006]、[Gaitsgory 2001]、[Frenkel and Gaitsgory 2006a]（均见本书 p.359 引用）。]}

% [Tutor Note: p.359 背景。Bezrukavnikov 的等价 (10.4-18) 是“仿射 Hecke 代数 $K$-理论实现”的范畴化：左边是仿射旗簇上扭变 D-模，右边是加厚 Steinberg 簇上的拟凝聚层。它与主猜想的关系如同“局部—整体”：中心层（来自 $D^b(\mathcal{D}^{\mathrm{Fl}}_{\kappa_c}\text{-}\mathrm{mod})^I$）在左边给出 Hecke 作用，而在右边对应 Steinberg 簇上的自然层。]
这个讨论表明范畴 $D^b(\hat{\fg}_{\kappa_c}\text{-}\mathrm{mod}_{\mathrm{nilp}})$ 与仿射旗簇上扭变 D-模范畴之间的密切联系，它类似于 $\hat{\fg}_{\kappa_c}\text{-}\mathrm{mod}_{\mathrm{reg}}$ 与仿射 Grassmann 流形上扭变 D-模范畴之间的联系（后者在 10.3.6 节讨论过）。联系 $D^b(\hat{\fg}_{\kappa_c}\text{-}\mathrm{mod}_{\mathrm{nilp}})$ 与 $D^b(\mathcal{D}^{\mathrm{Fl}}_{\kappa_c}\text{-}\mathrm{mod})$ 的更精确猜想在 [Frenkel and Gaitsgory 2006c] 中给出（见引言与第 6 节），读者可参看该文了解详情。这个猜想可以看作 猜想 10.3.9 对幂零 opers 的类比。正如 [Frenkel and Gaitsgory 2006c] 所解释，这个猜想由 [Arkhipov and Bezrukavnikov 2002; Arkhipov et al. 2004] 的结果支持（另见 [Bezrukavnikov 2006]）。这些结果与猜想合起来，为等价 (10.4-13) 提供了额外的证据。

\subsection{\texorpdfstring{$\fg=\fsl_2$}{g=sl2} 的情形}

\begin{keypoints}
    \item \textbf{核心思想：} 在 $\fg=\fsl_2$ 时把主猜想的各个要素显式化：$\mathrm{MOp}^0_{PGL_2}$ 的两个分支、Miura 变换的纤维、Wakimoto 模与扩张。
    \item \textbf{数学工具：} 非约化 Springer 纤维 $\Spec\C[\epsilon]/(\epsilon^2)$；Riccati 方程与递推公式 (10.4-19)、(10.4-20)；Wakimoto 模 $W_{u(t)}$；$\operatorname{Ext}^1$ 与 $\PP^1$。
    \item \textbf{关键结论：} Miura 变换在正则 oper 上的纤维是 $\PP^1$（对应 $\operatorname{Ext}^1$ 的射影化），在非零留数处是双点；主猜想在此情形给出 $\hat\fsl_2\text{-}\mathrm{mod}_\chi^{I^0}\simeq\mathrm{QCoh}(\Spec h_0)$。
    \item \textbf{记号提醒：} $u_{\mathrm{red}}(t)$ 是 $u(t)\bmod\epsilon$ 的解；$u_a(t)$ 是 $\C[\epsilon]/(\epsilon^2)$ 上的单参数解族；$\mathbb{V}_0(\chi)$ 与 $\mathbb{V}_{-2}(\chi)$ 是两个不可约模。
\end{keypoints}

% [Tutor Note: p.360--364。本小节是全节的“工作例子”：所有抽象陈述在 $\fsl_2$ 时都变成可算的。请抓住三条线：(1) $\widetilde{\widetilde{\mathrm{n}}}$ 的两个分支（$\PP^1$ 与加倍仿射线）；(2) Miura 变换的 Riccati 方程递推与自由度（正则情形 1 个自由参数 $\leftrightarrow\PP^1$，非零留数情形 0 个自由参数 $\leftrightarrow$ 点，但 $\C[\epsilon]/(\epsilon^2)$ 给出双点）；(3) Wakimoto 模作为扩张。]
本节用 $\fg=\fsl_2$ 的情形来说明 猜想 10.4.4，从而更具体地解释该猜想的关键要素。

首先讨论概形 $\widetilde{\mathrm{MOp}}{}^0_{PGL_2}$ 的结构。按定义，
\[
    \widetilde{\mathrm{MOp}}{}^0_{PGL_2}
    =\Op^{\mathrm{nilp}}_{PGL_2}
    \times_{\mathrm{n}/PGL_2}
    \widetilde{\widetilde{\mathrm{n}}}/PGL_2.
\]
簇 $\widetilde{\widetilde{\mathrm{n}}}$ 映入 $\mathrm{n}=\C$，它在 $u\in\mathrm{n}$ 上的纤维是 $u$ 的非约化 Springer 纤维，即 $\fsl_2$ 中含 $u$ 的 Borel 子代数构成的簇。

有两种可能。若 $u=0$，纤维是 $\PP^1$，即 $PGL_2$ 的旗簇。若 $u\neq0$，它对应正则幂零元，因此含 $u$ 的 Borel 子代数唯一。因此约化 Springer 纤维是单点。但非约化 Springer 纤维是双点 $\Spec\C[\epsilon]/(\epsilon^2)=\Spec h_0$。

为看出这一点，我们证明 Grothendieck alteration $\widetilde{\fsl_2}$ 有非平凡的 $\C[\epsilon]/(\epsilon^2)$-点，它们投影到 $\fsl_2$ 的正则幂零元。不妨取此元素由矩阵
\[
    e=\begin{pmatrix}0&1\\0&0\end{pmatrix}
\]
表示。则 $e$ 上方 $\widetilde{\fsl_2}$ 的一个 $\C[\epsilon]/(\epsilon^2)$-点，是 $\fsl_2[\epsilon]/(\epsilon^2)$ 中含 $e$（模 $\epsilon$ 约化后）的 Borel 子代数。现在注意：$\fsl_2[\epsilon]/(\epsilon^2)$ 中上三角矩阵的李代数 $\mathfrak{b}[\epsilon]/(\epsilon^2)$ 含元素
\[
    \begin{pmatrix}1&0\\-a\epsilon&1\end{pmatrix}
    e
    \begin{pmatrix}1&0\\a\epsilon&1\end{pmatrix}
    =
    \begin{pmatrix}a\epsilon&1\\0&-a\epsilon\end{pmatrix}
\]
对任意 $a\in\C$。因此
\[
    e\in
    \begin{pmatrix}1&0\\-a\epsilon&1\end{pmatrix}
    \mathfrak{b}[\epsilon]/(\epsilon^2)
    \begin{pmatrix}1&0\\a\epsilon&1\end{pmatrix},
    \qquad a\in\C.
\]
于是我们找到 Springer 纤维在 $e$ 上方的非平凡 $\C[\epsilon]/(\epsilon^2)$-点。

因此得到：簇 $\widetilde{\widetilde{\mathrm{n}}}$ 有两个分支，其一是 $\PP^1$，其二是加倍的仿射线 $\widetilde{\mathbb{A}}^1$。这两个分支对应 $\fsl_2$ 的外尔群的两个元素，它们标记 $\fsl_2$ 中两个 Borel 子代数的可能相对位置。直线上的点对应 Borel 子代数与含 $\mathrm{n}$ 的固定 Borel 子代数 $\mathfrak{b}$ 重合的情形，而 $\PP^1$ 的一般点对应与 $\mathfrak{b}$ 处于一般相对位置的 Borel 子代数。这两个分支在一点处“粘合”在一起。于是我们发现 $\widetilde{\widetilde{\mathrm{n}}}$ 是奇异、可约、非约化的拟射影簇。

这让我们对 $\widetilde{\mathrm{MOp}}{}^0_{PGL_2}$ 的样子有了粗略概念。它纤维化于幂零 $PGL_2$-opers 的空间 $\Op^{\mathrm{nilp}}_{PGL_2}$ 之上。取圆盘上的坐标 $t$，用它把 $\Op^{\mathrm{nilp}}_{PGL_2}$ 等同于二阶算子
\[
    \partial_t^2-v(t),\qquad v(t)=\sum_{n\ge-1}v_nt^n
\]
的空间。容易看出留数映射 $\Op^{\mathrm{nilp}}_{PGL_2}\to\mathrm{n}=\C$ 由公式 $v(t)\mapsto v_{-1}$ 给出。因此，若把 $\Op^{\mathrm{nilp}}_{PGL_2}$ 分解为坐标为 $v_{-1}$ 的直线与坐标为 $v_n$（$n\ge0$）的无穷维仿射空间的笛卡尔积，就得到：$\widetilde{\mathrm{MOp}}{}^0_{PGL_2}$ 同构于 $\widetilde{\widetilde{\mathrm{n}}}$ 与一个无穷维仿射空间的笛卡尔积。

于是 $\widetilde{\mathrm{MOp}}{}^0_{PGL_2}$ 有两个分支 $\widetilde{\mathrm{MOp}}{}^{0,1}_{PGL_2}$ 与 $\widetilde{\mathrm{MOp}}{}^{0,w_0}_{PGL_2}$，对应 $\fsl_2$ 的外尔群的两个元素 $1$ 与 $w_0$。

猜想 10.4.4 断言：概形 $\widetilde{\mathrm{MOp}}{}^0_{PGL_2}$ 上拟凝聚层的导出范畴，等价于等变导出范畴
\[
    D^b\bigl(\hat\fsl_2\text{-}\mathrm{mod}_{\mathrm{nilp}}\bigr)^{I^0}.
\]
下面考察这些范畴中对象的一些例子。

% [Tutor Note: p.361 易错处。$\widetilde{\mathrm{MOp}}{}^0_{PGL_2}\cong\widetilde{\widetilde{\mathrm{n}}}\times(\text{无穷维仿射空间})$：把 $\Op^{\mathrm{nilp}}_{PGL_2}$ 里除 $v_{-1}$ 外的高阶系数看成“自由方向”。这个分解使得所有非平凡几何都集中在 $v_{-1}$（即留数）方向上，这正是 $\fsl_2$ 例子可算的根源。]
首先考虑 $\widetilde{\mathrm{MOp}}{}^0_{PGL_2}$ 一点处的摩天楼层，它等同于 $\mathrm{MOp}^0_{PGL_2}$ 的一点。簇 $\mathrm{MOp}^0_{PGL_2}$ 是两个层片之并：
\[
    \mathrm{MOp}^0_{PGL_2}
    =\mathrm{MOp}^{0,1}_{PGL_2}
    \bigsqcup
    \mathrm{MOp}^{0,w_0}_{PGL_2}.
\]
由定理 9.3.6，
\[
    \mathrm{MOp}^{0,w}_{PGL_2}
    \simeq\mathrm{Conn}(\Omega^{-\rho})^{w(\rho)-\rho}_D
\]
（这里像以前一样从 $\Omega^{\check\rho}$ 换到 $\Omega^{-\rho}$）。

对 $w=1$，层片 $\mathrm{MOp}^{0,1}_{PGL_2}$ 纤维化于正则 opers 的子空间 $\Op_{PGL_2}(D)\subset\Op^{\mathrm{nilp}}_{PGL_2}$（即留数为零、$v_{-1}=0$ 的那些）之上。相应的联络空间 $\mathrm{Conn}(\Omega^{-\rho})^{w(\rho)-\rho}_D$ 是 $D$ 上线丛 $\Omega^{-1/2}$ 上正则联络的空间。用坐标 $t$，把这样的联络表示为算子
\[
    \overline{\nabla}=\partial_t+u(t),\qquad u(t)=\sum_{n\ge0}u_nt^n.
\]
映射 $\mathrm{MOp}^{0,1}_{PGL_2}\to\Op_{PGL_2}(D)$ 就是 Miura 变换
\[
    \partial_t+u(t)\longmapsto
    (\partial_t-u(t))(\partial_t+u(t))=\partial_t^2-v(t),
\]
其中
\begin{equation}
    v(t)=u(t)^2-\partial_tu(t).
    \tag{10.4-19}
\end{equation}
这些公式此前已在 8.2.2 节得到，只是现在把 $u(t)$ 重标度了 $2$ 倍以简化计算。

Miura 变换 $\mathrm{Conn}(\Omega^{-1/2})_D\to\Op_{PGL_2}(D)$ 在 $v(t)\in\C[[t]]\simeq\Op_{PGL_2}(D)$ 上的纤维，由 Riccati 方程 (10.4-19) 的一切解 $u(t)\in\C[[t]]$ 组成。这个方程等价于关于 $u(t)$ 的系数 $u_n$ 的无穷代数方程组，可以递推求解。这些方程是
\[
    nu_n=-v_{n-1}+\sum_{i+j=n-1}u_iu_j,\qquad n\ge1.
\]
于是我们看到，对 $u_0$ 的每个选取（$u_0$ 可以是任意复数），方程唯一确定所有系数 $u_n$（$n\ge1$）。因此这个映射的纤维同构于仿射直线，正如所期待。

对上述方程的每个解 $u(t)\in\C[[t]]$，有推论 6.1.5 定义的 Wakimoto 模 $W_{u(t)}$。根据定理 8.3.3，中心 $Z(\hat\fsl_2)$ 通过对应于 $PGL_2$-oper $\chi$（对应射影联络 $\partial_t^2-v(t)$）的中心特征标作用在 $W_{u(t)}$ 上。于是得到范畴 $\hat\fsl_2\text{-}\mathrm{mod}_\chi$ 中由仿射直线参数化的一族 Wakimoto 模。这些就是对应支撑在 $\mathrm{MOp}^{0,1}_{PGL_2}$ 各点上的摩天楼层的对象。

% [Tutor Note: p.362 易错处。Riccati 递推 $nu_n=-v_{n-1}+\sum_{i+j=n-1}u_iu_j$（$n\ge1$）中，$u_0$ 是自由参数，其余 $u_n$ 由它唯一决定；这正是“Miura 变换在正则 oper 上的纤维 ≅ 仿射直线”的代数体现。注意 $v_{n-1}$ 前的负号来自 $v=u^2-\partial_tu$。]
现在考虑另一层片
\[
    \mathrm{MOp}^{0,w_0}_{PGL_2}
    \simeq\mathrm{Conn}(\Omega^{-1/2})^{-1}_D
\]
的点。$\Omega^{-1/2}$ 上相应的联络在 $0\in D$ 处有正则奇点，留数为 $-1$（对应 $-2\rho$）。用我们的坐标 $t$，它们对应算子
\[
    \overline{\nabla}=\partial_t+u(t),\qquad
    u(t)=\sum_{n\ge-1}u_nt^n,\quad u_{-1}=-1.
\]
Miura 变换 $\mathrm{Conn}(\Omega^{-1/2})^{-1}_D\to\Op^{\mathrm{nilp}}_{PGL_2}$ 仍由公式 (10.4-19) 给出。下面确定该映射在对应于射影联络 $\partial_t^2-v(t)$ 的幂零 oper 上的纤维。得到关于 $u(t)$ 的系数 $u_n$ 的如下方程组
\begin{equation}
    (n+2)u_n=-v_{n-1}
    +\sum_{i+j=n-1;\ i,j\ge0}u_iu_j,
    \qquad n\ge0
    \tag{10.4-20}
\end{equation}
（这里用了 $u_{-1}=-1$）。这个方程组可以递推求解，但现在 $u_0$ 由第一个方程确定，因此不像前一情形那样有自由参数。于是我们发现映射 $\mathrm{Conn}(\Omega^{-1/2})^{-1}_D\to\Op^{\mathrm{nilp}}_{PGL_2}$ 是同构，正如所期待。

若 $v_{-1}=0$，即我们的 oper $\chi=\partial_t^2-v(t)$ 是正则的，则在 Miura 变换在 $\chi$ 上的纤维中多得到一个点。因此我们发现，正则 oper $\chi$ 上方纤维的点集确实与射影直线的点集双射。Wakimoto 模 $W_{u(t)}$（其中 $u(t)$ 如上）是范畴 $\hat\fsl_2\text{-}\mathrm{mod}_\chi^{I^0}$ 中对应于由 $u(t)$ 表示的额外点的对象。于是对每个正则 oper $\chi$，得到范畴 $\hat\fsl_2\text{-}\mathrm{mod}_\chi^{I^0}$ 中由 $\PP^1$ 参数化的一族 Wakimoto 模，正如所期待。

现在设 $v_{-1}\neq0$，于是 oper $\chi$ 的留数非零，从而它有非平凡幂幺单值性。则 Miura 变换在 $\chi$ 上的集合论纤维由单点组成。我们有一个 Wakimoto 模 $W_{u(t)}$，其中 $u(t)=-1/t+\cdots$，对应该点。然而 $\widetilde{\mathrm{MOp}}{}^0_{PGL_2}$ 的纤维是双点。从 Miura 变换的角度能看出这个加倍效应吗？

答案是可以。考虑方程 (10.4-19)，其中 $v_{-1}\neq0$，$u(t)=\sum_{n\in\Z}u_nt^n$，$u_n\in\C[\epsilon]/(\epsilon^2)$。容易看出，满足 $u(t)\bmod\epsilon$ 为上述解 $u_{\mathrm{red}}(t)=-1/t+\sum_{n\ge0}u_{\mathrm{red},n}t^n$ 的一般解具有形式
\[
    u_a(t)=u_{\mathrm{red}}(t)+\frac{a\epsilon}{t^2}
    +a\epsilon\sum_{n\ge-1}\delta_nt^n,
\]
其中 $a$ 是任意复数，$\delta_n\in\C$ 由 (10.4-19) 唯一确定。于是得到：对上述形式的 $\chi$（$v_{-1}\neq0$），Miura 变换 $\mathrm{Conn}(\Omega^{-\rho})_{D^\times}\to\Op^{\mathrm{nilp}}_{PGL_2}$ 的概形论纤维是双点，就像 $\widetilde{\mathrm{MOp}}{}^0_{PGL_2}$ 的纤维一样。

这提示：一般地，Miura 变换在幂零 oper $\chi\in\Op^{\mathrm{nilp}}_{{}^L G}$ 上的概形论纤维，也同构于非约化 Springer 纤维 $\operatorname{Sp}_{\Res_{\mathcal{F}}(\chi)}$。在 [Frenkel and Gaitsgory 2006c] 中已证明确实如此。

% [Tutor Note: p.363 背景。这里用 $\C[\epsilon]/(\epsilon^2)$ 上的解族 $u_a(t)$ 解释“双点”：$u_a(t)$ 的线性化参数 $a$ 对应一维自扩张空间。也就是说，非约化 Springer 纤维的“加厚方向”在代数一侧对应 Wakimoto 模的一维自扩张。这种“几何加倍 ↔ 模自扩张”的对应是全节最具体的证据之一。]
上述 (10.4-19) 在 $\C[\epsilon]/(\epsilon^2)$ 上的单参数解族 $u_a(t)$（$a\in\C$）的存在，意味着 Wakimoto 模 $W_{u_{\mathrm{red}}(t)}$ 有一维的自扩张空间。回忆这个模实现在向量空间 $M_{\fsl_2}$ 中，即由 $a_n,a_n^*$（$n\in\Z$）生成的外尔代数的 Fock 模。这个扩张实现在 $M_{\fsl_2}\otimes\C[\epsilon]/(\epsilon^2)$ 中。李代数 $\hat\fsl_2$ 的作用由与以前相同的公式给出（见 6.1.2 节），只是现在 $b_n$（$n\in\Z$）作为 $\C[\epsilon]/(\epsilon^2)$ 上的线性算子作用，对应于 $u_a(t)$ 的第 $n$ 个系数 $u_{a,n}$。

事实上，由定理 10.4.5（或由显式计算）可推出：对应于留数非零的 $\chi\in\Op^{\mathrm{nilp}}_{PGL_2}$ 的范畴 $\hat\fsl_2\text{-}\mathrm{mod}_\chi^{I^0}$ 有唯一的不可约对象，即 $W_{u_{\mathrm{red}}(t)}$。我们指出，这个模同构于 Verma 模 $\mathbb{M}_{-2}$ 模去对应于 $\chi$ 的中心特征标所得的商。这个模有上述公式描述的非平凡自扩张，且范畴 $\hat\fsl_2\text{-}\mathrm{mod}_\chi^{I^0}$ 等价于范畴 $\mathrm{QCoh}(\Spec h_0)$，其中 $h_0=\C[\epsilon]/(\epsilon^2)$。

现在考虑范畴 $\hat\fsl_2\text{-}\mathrm{mod}_\chi^{I^0}$，其中 $\chi$ 是对应于算子 $\partial_t^2-v(t)$ 的 $PGL_2$-oper，且在 $t=0$ 处正则（从而留数 $v_{-1}=0$）。我们构造了该范畴的一族对象 $W_{u(t)}$，其中 $u(t)$ 是 Riccati 方程 (10.4-19) 的解，它们与 $\PP^1$ 的点一一对应（除一个之外，$u(t)$ 都在 $t=0$ 处正则，最后那个具有形式 $u(t)=-1/t+\cdots$）。这些 Wakimoto 模的结构如何？不难直接看出，它们每个都是同两个不可约 $\hat\fsl_2$-模的扩张，而这两个模事实上是范畴 $\hat\fsl_2\text{-}\mathrm{mod}_\chi^{I^0}$ 仅有的不可约对象。第一个是 10.3.3 节引入的模 $\mathbb{V}_0(\chi)$，第二个由 $\mathbb{V}_0(\chi)$ 经 $\hat\fsl_2$ 的对合扭变得到，该对合在 Kac--Moody 生成元上（见附录 A.5）由公式
\[
    e_0\leftrightarrow e_1,\qquad
    h_0\leftrightarrow h_1,\qquad
    f_0\leftrightarrow f_1
\]
定义。我们把这个模记作 $\mathbb{V}_{-2}(\chi)$。$\mathbb{V}_0(\chi)$ 由被“极大紧”李子代数 $\mathfrak{k}_0=\fsl_2[[t]]$ 零化的向量生成，而 $\mathbb{V}_{-2}(\chi)$ 由被李子代数 $\mathfrak{k}'_0$ 零化的向量生成，$\mathfrak{k}'_0$ 是上述对合作用于 $\fsl_2[[t]]$ 得到的。这两个李子代数 $\mathfrak{k}_0$ 与 $\mathfrak{k}'_0$，是 $\hat\fsl_2$ 在 $SL_2((t))$ 共轭下的两个“极大紧”李子代数。特别地，$\mathbb{V}_{-2}(\chi)$ 的生成向量被李代数 $\widehat{\fn}_+$ 零化，且 Cartan 子代数 $\fh$ 按权 $-2$（对应 $-\alpha$）作用在它上面。因此 $\mathbb{V}_{-2}(\chi)$ 是 Verma 模 $\mathbb{M}_{-2}$ 的商。

我们有二维的扩张空间 $\operatorname{Ext}^1(\mathbb{V}_{-2}(\chi),\mathbb{V}_0(\chi))$。由于对应于 $\operatorname{Ext}^1(\mathbb{V}_{-2}(\chi),\mathbb{V}_0(\chi))$ 中成比例向量的扩张作为 $\hat\fsl_2$-模同构，我们得到扩张的同构类由 $\PP^1$ 的点参数化。它们恰好由我们的 Wakimoto 模 $W_{u(t)}$ 表示。

% [Tutor Note: p.364。这是本节最后一个具体计算，也是连接 p.375 的接点：$\operatorname{Ext}^1(\mathbb{V}_{-2}(\chi),\mathbb{V}_0(\chi))$ 是二维的，其射影化就是 $\PP^1$，即 Miura 变换在正则 oper 处的纤维。注意这里只用到“二维扩张空间”，其射影化 $\PP^1$ 与 Miura 纤维的 $\PP^1$ 的等同需要下一句（在 PDF p.375，由下一单元接译）来完成。]
% 注意：本句在 PDF p.375（印刷 p.365）续完；下一单元 U17 接译“acquires a useful interpretation as the projectivization of the space of extensions of the two irreducible objects of the category...”。
因此，我们的 $\PP^1$，即 Miura 变换在正则 oper 处的纤维，
